\documentclass[11pt]{article}

\usepackage[margin=1in]{geometry}
% Shared preamble for the urcu-txn paper series.
% Included by each paper's main.tex AFTER \documentclass.
%
% Deliberately targets what arXiv's AutoTeX accepts:
%   - pdflatex (arXiv default is TeX Live 2025, matching our local install)
%   - natbib + bibtex, with the generated .bbl committed alongside the source
%     (biblatex .bbl is version-locked to the biblatex/biber pair; natbib's
%     is not, so natbib is the lower-risk path even though arXiv now runs
%     both backends)
%   - pgfplots reading CSV directly, so plot data ships inside the submission
%     and no external plotting toolchain is needed to rebuild a figure

\usepackage[T1]{fontenc}
\usepackage[utf8]{inputenc}
\usepackage{lmodern}
\usepackage{microtype}
\usepackage{xspace}
\usepackage{xcolor}

% This series is dense with long, unbreakable C identifiers
% (cds_list_for_each_entry_reverse, -DURCU_TXN_SW_EXCL_VALIDATE), which TeX
% cannot hyphenate and will happily shove into the margin. emergencystretch
% lets it loosen interword spacing as a last resort instead.
%
% Why not hyphenate typewriter text (\usepackage[htt]{hyphenat}): a hyphen
% inside a C identifier is ambiguous -- the reader cannot tell whether it is
% part of the name. A slightly loose line is a much cheaper defect than an
% identifier the reader has to guess at, in a paper whose point is enablement.
\setlength{\emergencystretch}{3em}

\usepackage{amsmath}
\usepackage{amssymb}
\usepackage{amsthm}

\usepackage{booktabs}
\usepackage{multirow}
\usepackage{graphicx}
\usepackage{subcaption}
\usepackage{siunitx}

\usepackage{tikz}
% positioning: `above=2mm of foo`. Without it TikZ parses "of" as a math
% operator and dies with an opaque PGF Math Error rather than a missing-library
% message, so load it up front.
\usetikzlibrary{positioning,arrows.meta,calc,fit,backgrounds}
\usepackage{pgfplots}
\pgfplotsset{compat=1.18}
\usepgfplotslibrary{groupplots}

\usepackage[ruled,vlined,linesnumbered]{algorithm2e}
\usepackage{listings}
\usepackage{etoolbox}

\usepackage[numbers,sort&compress]{natbib}

% hyperref late, cleveref last.
\usepackage[colorlinks=true,linkcolor=black,citecolor=black,urlcolor=blue]{hyperref}
\usepackage{cleveref}

% ---------------------------------------------------------------------------
% C listings style, for the API and pseudocode excerpts. Enablement is the
% point of a defensive publication, so code excerpts are first-class content
% and must stay legible in print.
% ---------------------------------------------------------------------------
\lstdefinestyle{c}{
  language=C,
  basicstyle=\ttfamily\small,
  keywordstyle=\bfseries,
  commentstyle=\itshape\color{gray!70!black},
  stringstyle=\color{black},
  numbers=none,
  frame=none,
  columns=fullflexible,
  keepspaces=true,
  showstringspaces=false,
  breaklines=true,
  captionpos=b,
  morekeywords={uint64_t,uintptr_t,bool,struct,static,inline,_Atomic},
}
\lstset{style=c}

% ---------------------------------------------------------------------------
% Listings must not split across a page boundary. The excerpts in this series
% are short (a handful of lines each) and are read as single units -- a
% __list_del() whose two stores land on different pages defeats the entire
% point of showing them together. listings has no option for this, so wrap
% every listing in an unbreakable box. If a listing ever grows longer than a
% page this will overfull rather than break: that is the correct failure, since
% such a listing wants to be a float with a caption instead.
% ---------------------------------------------------------------------------
\BeforeBeginEnvironment{lstlisting}{\par\noindent\begin{minipage}{\linewidth}}
\AfterEndEnvironment{lstlisting}{\end{minipage}\par}

\theoremstyle{definition}
\newtheorem{definition}{Definition}
\theoremstyle{plain}
\newtheorem{theorem}{Theorem}
\newtheorem{lemma}{Lemma}
\newtheorem{invariant}{Invariant}
\theoremstyle{remark}
\newtheorem{property}{Property}

% Shared macros for the urcu-txn paper series.
% Included by each paper's main.tex via % Shared macros for the urcu-txn paper series.
% Included by each paper's main.tex via % Shared macros for the urcu-txn paper series.
% Included by each paper's main.tex via \input{../common/macros}.
%
% Keep this class-agnostic: nothing here may assume `article`, so that a paper
% can be retargeted to acmart/IEEEtran by changing only its own preamble.

% ---------------------------------------------------------------------------
% Terminology.
%
% These exist so the series says the same thing the same way every time, and so
% that a terminology decision is changed in one place. The guarantee wording is
% load-bearing: this engine is NOT STM and is NOT serializable. See
% \pseudotxn below and the note in README.md.
% ---------------------------------------------------------------------------
\newcommand{\pseudotxn}{pseudo-transaction\xspace}
\newcommand{\pseudotxns}{pseudo-transactions\xspace}
\newcommand{\mcas}{MCAS\xspace}
\newcommand{\kcas}{$k$-CAS\xspace}
\newcommand{\rcu}{RCU\xspace}
\newcommand{\QSBR}{QSBR\xspace}
\newcommand{\urcu}{userspace RCU\xspace}
\newcommand{\sw}{single-writer\xspace}
\newcommand{\mw}{multi-writer\xspace}
\newcommand{\fliplatch}{flip-latch\xspace}
\newcommand{\ryw}{read-your-own-writes\xspace}

% Engine identifiers.
\newcommand{\urcutxn}{\texttt{urcu-txn}\xspace}
\newcommand{\rcumcas}{\texttt{rcu-mcas}\xspace}
\newcommand{\rcutxnsw}{\texttt{rcu-txn-sw}\xspace}

% ---------------------------------------------------------------------------
% Code and identifiers.
% ---------------------------------------------------------------------------
% \code offers TeX no place to break, and this series is dense with identifiers
% long enough to overflow a line on their own (cds_list_for_each_entry_reverse,
% -DURCU_TXN_SW_EXCL_VALIDATE). Two of them in one sentence is enough to push
% text into the margin no matter how the surrounding prose is stretched.
%
% Since every such identifier already writes its underscores as \_, redefining
% \_ locally gives TeX a legal breakpoint after each one -- silently, with NO
% hyphen. That last part is the whole reason not to reach for hyphenat's [htt]
% option instead: a hyphen inside a C identifier is ambiguous, and the reader
% cannot tell whether it belongs to the name. For a paper whose purpose is
% enablement, an identifier someone has to guess at is a real defect; a slightly
% loose line is not.
%
% Scoped to \texttt's own group, so \_ keeps its normal meaning everywhere else.
%
% \let, not \renewcommand, and \DeclareRobustCommand, not \newcommand: \code is
% used inside \caption, which is a MOVING argument (written to the .aux and
% re-read). A fragile \code expands there and the \renewcommand form dies with
% "Argument of \caption@ydblarg has an extra }" -- a fatal error whose message
% names neither \code nor the caption it came from.
\newcommand{\codeunderscore}{\textunderscore\allowbreak}
\DeclareRobustCommand{\code}[1]{\texttt{\let\_\codeunderscore #1}}
\DeclareRobustCommand{\fn}[1]{\code{#1()}}

% A record in a transaction's write set: {slot, old, new}.
\newcommand{\record}[3]{$\{#1,\,#2 \rightarrow #3\}$}
% An edit on one slot, old to new.
\newcommand{\edit}[2]{$#1 \rightarrow #2$}

% Transaction status values.
\newcommand{\UNDECIDED}{\textsc{undecided}\xspace}
\newcommand{\SUCCEEDED}{\textsc{succeeded}\xspace}
\newcommand{\FAILED}{\textsc{failed}\xspace}
\newcommand{\ABORT}{\textsc{abort}\xspace}

% ---------------------------------------------------------------------------
% Editorial markers. Define \draftmode in the preamble to make these visible;
% they vanish in a submission build. Never ship with \draftmode set.
% ---------------------------------------------------------------------------
\newif\ifdraftmode
\draftmodefalse

% These use \color/\bfseries SWITCHES inside a group rather than
% \textcolor{}{}/\textbf{}, whose arguments are short: an editorial note
% spanning a blank line would otherwise die with "Paragraph ended before
% \@textcolor was complete". Notes routinely span paragraphs, so the switch
% form is the only one that survives real use.
\newcommand{\marker}[3]{\ifdraftmode{\color{#1}\bfseries[#2: #3]}\fi}

\newcommand{\TODO}[1]{\marker{red}{TODO}{#1}}
\newcommand{\NOTE}[1]{\marker{blue}{NOTE}{#1}}
% Claims needing a citation before submission.
\newcommand{\NEEDCITE}[1]{\marker{orange}{CITE}{#1}}
% Numbers currently sourced from source comments rather than a reproducible run.
\newcommand{\UNVERIFIED}[1]{\marker{purple}{UNVERIFIED}{#1}}
.
%
% Keep this class-agnostic: nothing here may assume `article`, so that a paper
% can be retargeted to acmart/IEEEtran by changing only its own preamble.

% ---------------------------------------------------------------------------
% Terminology.
%
% These exist so the series says the same thing the same way every time, and so
% that a terminology decision is changed in one place. The guarantee wording is
% load-bearing: this engine is NOT STM and is NOT serializable. See
% \pseudotxn below and the note in README.md.
% ---------------------------------------------------------------------------
\newcommand{\pseudotxn}{pseudo-transaction\xspace}
\newcommand{\pseudotxns}{pseudo-transactions\xspace}
\newcommand{\mcas}{MCAS\xspace}
\newcommand{\kcas}{$k$-CAS\xspace}
\newcommand{\rcu}{RCU\xspace}
\newcommand{\QSBR}{QSBR\xspace}
\newcommand{\urcu}{userspace RCU\xspace}
\newcommand{\sw}{single-writer\xspace}
\newcommand{\mw}{multi-writer\xspace}
\newcommand{\fliplatch}{flip-latch\xspace}
\newcommand{\ryw}{read-your-own-writes\xspace}

% Engine identifiers.
\newcommand{\urcutxn}{\texttt{urcu-txn}\xspace}
\newcommand{\rcumcas}{\texttt{rcu-mcas}\xspace}
\newcommand{\rcutxnsw}{\texttt{rcu-txn-sw}\xspace}

% ---------------------------------------------------------------------------
% Code and identifiers.
% ---------------------------------------------------------------------------
% \code offers TeX no place to break, and this series is dense with identifiers
% long enough to overflow a line on their own (cds_list_for_each_entry_reverse,
% -DURCU_TXN_SW_EXCL_VALIDATE). Two of them in one sentence is enough to push
% text into the margin no matter how the surrounding prose is stretched.
%
% Since every such identifier already writes its underscores as \_, redefining
% \_ locally gives TeX a legal breakpoint after each one -- silently, with NO
% hyphen. That last part is the whole reason not to reach for hyphenat's [htt]
% option instead: a hyphen inside a C identifier is ambiguous, and the reader
% cannot tell whether it belongs to the name. For a paper whose purpose is
% enablement, an identifier someone has to guess at is a real defect; a slightly
% loose line is not.
%
% Scoped to \texttt's own group, so \_ keeps its normal meaning everywhere else.
%
% \let, not \renewcommand, and \DeclareRobustCommand, not \newcommand: \code is
% used inside \caption, which is a MOVING argument (written to the .aux and
% re-read). A fragile \code expands there and the \renewcommand form dies with
% "Argument of \caption@ydblarg has an extra }" -- a fatal error whose message
% names neither \code nor the caption it came from.
\newcommand{\codeunderscore}{\textunderscore\allowbreak}
\DeclareRobustCommand{\code}[1]{\texttt{\let\_\codeunderscore #1}}
\DeclareRobustCommand{\fn}[1]{\code{#1()}}

% A record in a transaction's write set: {slot, old, new}.
\newcommand{\record}[3]{$\{#1,\,#2 \rightarrow #3\}$}
% An edit on one slot, old to new.
\newcommand{\edit}[2]{$#1 \rightarrow #2$}

% Transaction status values.
\newcommand{\UNDECIDED}{\textsc{undecided}\xspace}
\newcommand{\SUCCEEDED}{\textsc{succeeded}\xspace}
\newcommand{\FAILED}{\textsc{failed}\xspace}
\newcommand{\ABORT}{\textsc{abort}\xspace}

% ---------------------------------------------------------------------------
% Editorial markers. Define \draftmode in the preamble to make these visible;
% they vanish in a submission build. Never ship with \draftmode set.
% ---------------------------------------------------------------------------
\newif\ifdraftmode
\draftmodefalse

% These use \color/\bfseries SWITCHES inside a group rather than
% \textcolor{}{}/\textbf{}, whose arguments are short: an editorial note
% spanning a blank line would otherwise die with "Paragraph ended before
% \@textcolor was complete". Notes routinely span paragraphs, so the switch
% form is the only one that survives real use.
\newcommand{\marker}[3]{\ifdraftmode{\color{#1}\bfseries[#2: #3]}\fi}

\newcommand{\TODO}[1]{\marker{red}{TODO}{#1}}
\newcommand{\NOTE}[1]{\marker{blue}{NOTE}{#1}}
% Claims needing a citation before submission.
\newcommand{\NEEDCITE}[1]{\marker{orange}{CITE}{#1}}
% Numbers currently sourced from source comments rather than a reproducible run.
\newcommand{\UNVERIFIED}[1]{\marker{purple}{UNVERIFIED}{#1}}
.
%
% Keep this class-agnostic: nothing here may assume `article`, so that a paper
% can be retargeted to acmart/IEEEtran by changing only its own preamble.

% ---------------------------------------------------------------------------
% Terminology.
%
% These exist so the series says the same thing the same way every time, and so
% that a terminology decision is changed in one place. The guarantee wording is
% load-bearing: this engine is NOT STM and is NOT serializable. See
% \pseudotxn below and the note in README.md.
% ---------------------------------------------------------------------------
\newcommand{\pseudotxn}{pseudo-transaction\xspace}
\newcommand{\pseudotxns}{pseudo-transactions\xspace}
\newcommand{\mcas}{MCAS\xspace}
\newcommand{\kcas}{$k$-CAS\xspace}
\newcommand{\rcu}{RCU\xspace}
\newcommand{\QSBR}{QSBR\xspace}
\newcommand{\urcu}{userspace RCU\xspace}
\newcommand{\sw}{single-writer\xspace}
\newcommand{\mw}{multi-writer\xspace}
\newcommand{\fliplatch}{flip-latch\xspace}
\newcommand{\ryw}{read-your-own-writes\xspace}

% Engine identifiers.
\newcommand{\urcutxn}{\texttt{urcu-txn}\xspace}
\newcommand{\rcumcas}{\texttt{rcu-mcas}\xspace}
\newcommand{\rcutxnsw}{\texttt{rcu-txn-sw}\xspace}

% ---------------------------------------------------------------------------
% Code and identifiers.
% ---------------------------------------------------------------------------
% \code offers TeX no place to break, and this series is dense with identifiers
% long enough to overflow a line on their own (cds_list_for_each_entry_reverse,
% -DURCU_TXN_SW_EXCL_VALIDATE). Two of them in one sentence is enough to push
% text into the margin no matter how the surrounding prose is stretched.
%
% Since every such identifier already writes its underscores as \_, redefining
% \_ locally gives TeX a legal breakpoint after each one -- silently, with NO
% hyphen. That last part is the whole reason not to reach for hyphenat's [htt]
% option instead: a hyphen inside a C identifier is ambiguous, and the reader
% cannot tell whether it belongs to the name. For a paper whose purpose is
% enablement, an identifier someone has to guess at is a real defect; a slightly
% loose line is not.
%
% Scoped to \texttt's own group, so \_ keeps its normal meaning everywhere else.
%
% \let, not \renewcommand, and \DeclareRobustCommand, not \newcommand: \code is
% used inside \caption, which is a MOVING argument (written to the .aux and
% re-read). A fragile \code expands there and the \renewcommand form dies with
% "Argument of \caption@ydblarg has an extra }" -- a fatal error whose message
% names neither \code nor the caption it came from.
\newcommand{\codeunderscore}{\textunderscore\allowbreak}
\DeclareRobustCommand{\code}[1]{\texttt{\let\_\codeunderscore #1}}
\DeclareRobustCommand{\fn}[1]{\code{#1()}}

% A record in a transaction's write set: {slot, old, new}.
\newcommand{\record}[3]{$\{#1,\,#2 \rightarrow #3\}$}
% An edit on one slot, old to new.
\newcommand{\edit}[2]{$#1 \rightarrow #2$}

% Transaction status values.
\newcommand{\UNDECIDED}{\textsc{undecided}\xspace}
\newcommand{\SUCCEEDED}{\textsc{succeeded}\xspace}
\newcommand{\FAILED}{\textsc{failed}\xspace}
\newcommand{\ABORT}{\textsc{abort}\xspace}

% ---------------------------------------------------------------------------
% Editorial markers. Define \draftmode in the preamble to make these visible;
% they vanish in a submission build. Never ship with \draftmode set.
% ---------------------------------------------------------------------------
\newif\ifdraftmode
\draftmodefalse

% These use \color/\bfseries SWITCHES inside a group rather than
% \textcolor{}{}/\textbf{}, whose arguments are short: an editorial note
% spanning a blank line would otherwise die with "Paragraph ended before
% \@textcolor was complete". Notes routinely span paragraphs, so the switch
% form is the only one that survives real use.
\newcommand{\marker}[3]{\ifdraftmode{\color{#1}\bfseries[#2: #3]}\fi}

\newcommand{\TODO}[1]{\marker{red}{TODO}{#1}}
\newcommand{\NOTE}[1]{\marker{blue}{NOTE}{#1}}
% Claims needing a citation before submission.
\newcommand{\NEEDCITE}[1]{\marker{orange}{CITE}{#1}}
% Numbers currently sourced from source comments rather than a reproducible run.
\newcommand{\UNVERIFIED}[1]{\marker{purple}{UNVERIFIED}{#1}}


% Draft build: `make draft` drops draft.flag to surface TODO/CITE/NOTE markers.
% Submission builds leave it absent, so the markers vanish.
\IfFileExists{draft.flag}{\draftmodetrue}{}

\title{RCU Pseudo-Transactions and the Sole-Driver MCAS:\\
       A Bounded-Blocking Multi-Writer Facility for\\
       Atomic Multi-Slot Update}

\author{Mathieu Desnoyers\\ EfficiOS Inc.\\ \texttt{mathieu.desnoyers@efficios.com}}

\date{\today}

\begin{document}
\maketitle

\begin{abstract}
% ---------------------------------------------------------------------------
% FRAMING RULE FOR THE WHOLE SERIES -- do not soften this in editing.
% This paper refuses to overclaim the PROGRESS class the way the companion
% refused to overclaim the CONSISTENCY model. The facility is bounded-blocking,
% NOT lock-free, and that is a deliberate tail-latency-for-throughput trade.
% Keep "linearizable" scoped to the k-CAS write. NEVER "serializable"/"opacity".
% Overclaiming here is the single most damaging error available to this paper.
% ---------------------------------------------------------------------------
A companion paper described a Read-copy-update (\rcu) \pseudotxn---a frozen set of
\record{slot}{old}{new} records made visible by one commit step---and realized
it for a single writer with a \fliplatch that has no compare-and-swap and
requires writer exclusion. This paper drops the exclusion, and drops it per
\emph{slot}: one commit spans locked and unlocked slots alike, and a reader
cannot tell them apart. The sole-driver
multi-word compare-and-swap (\mcas) facility realizes the same model for concurrent
writers: a transaction is driven only by its owner, and a thread meeting
an in-flight proxy waits a bounded spin for the owner to settle,
then escalates. There is no helping, no stealing, and---most of all---no RDCSS
(restricted double-compare single-swap). A plain value-CAS
install is A-B-A-safe by construction: the re-plant hazard it prevents
needs a \emph{foreign} driver of one's own records, which sole-driver excludes.

The facility is \emph{bounded-blocking}: it waits on a foreign owner, so it is
not lock-free, and a FIFO-fair mutex fallback makes it starvation-free under a
fair scheduler. We argue this is the right target, not a shortfall, and the
argument is measured: a helping design was built and removed. Helping forced a
per-record install latch that made it bounded-blocking anyway---so lock-freedom
was never on the table---and it lost on throughput.

Per-slot atomicity is not operation atomicity, so a usable concurrent list needs
one more thing: logical deletion, a tombstone that update transactions treat as
freezing a node. It is updater-only
coordination---a would-be attacher validates it and aborts---so a reader gets a
frozen node for free without consulting the tombstone at all, which takes it off
the reader's path and leaves its placement free. This is a known
technique, disclosed as applied. With it and one guarded read, the same
bidirectional list and hash list become usable under concurrent writers.
\end{abstract}

% ===========================================================================
\section{Introduction}
\label{sec:intro}

Read-copy-update (\rcu) rests on a bargain: readers pay nothing, and in exchange a
writer gets exactly one atomic operation, a single pointer store bracketed by building
the new structure where readers cannot reach it and reclaiming the old one after
a grace period~\citep{mckenney1998rcu, desnoyers2012urcu}. A companion paper
gave the \rcu writer a \emph{second} atomic operation: a \pseudotxn, a frozen set of
\record{slot}{old}{new} records that a single commit step makes visible at once,
so one store moves a set of slots rather than one. It realized that model for a
single writer with a \fliplatch---a mechanism with no compare-and-swap, no
conflict detection, and no abort, whose price is that it is legal only under
writer mutual exclusion.

This paper drops the exclusion. The question it answers is what the same model
costs when several writers commit concurrently, each without a lock. The answer
is a facility, the sole-driver \mcas (multi-word compare-and-swap), and a claim
about what \emph{kind} of facility it should be.

Exclusion turns out to be droppable \emph{per slot} rather than wholesale. Every
record carries a \emph{kind}: a shared slot is planted with a compare-and-swap,
an exclusive slot---one the embedder holds a lock over---is parked with a plain
store that cannot fail, and one commit carries both against a single
linearization point (\cref{sec:kinds}). The generalization costs the reader
nothing, because both kinds resolve through the same status word and a reader
cannot tell them apart; and it costs an unmixed write set nothing either,
because the commit skips the machinery such a write set does not need.

\paragraph{The thesis.} The companion paper refused to overclaim the model's
\emph{consistency}: it is a \pseudotxn, not software transactional memory, and
we said what isolation it declines and why. This paper refuses to overclaim the
facility's \emph{progress}. It is \emph{bounded-blocking}---it waits, in bounded
steps, on a foreign owner---rather than lock-free, and that is a deliberate
choice of synchronization coarseness: coarse enough to stay simple, portable,
and cache-friendly, fine enough to stay parallel. Bounded-blocking trades a
tail-latency property for throughput knowingly; it is not a failed attempt to
reach lock-freedom. The paper states what it gives up and argues the trade. It
does not apologize for it.

The argument is not a preference stated after the fact. It is a falsification
trail. An earlier version of this engine had threads \emph{help} drive each
other's transactions, in the manner of practical \mcas
\citep{harris2002mcas, barnes1993method}. Helping made two threads co-install a
transaction's records, which admits an install-time A-B-A hazard---a value that
cycles away and back, which a plain compare-and-swap cannot detect---which forced
a per-record \code{FREE}/\code{BUSY}/\code{DONE} install latch to close. That latch
already made the helping engine bounded-blocking---so lock-freedom was never
actually on the table there---and it lost on throughput to the simpler design.
The cost is countable: the latch put \emph{three extra
compare-and-swaps on every record}---one to claim the install
(\code{FREE}\,$\to$\,\code{BUSY}), and two at settle, where a sole driver needs
only a release store. The decision itself went the same way, from
a single release store of the status word to a compare-and-swap, because an
owner that may be raced cannot assume it is the only writer of its own outcome.
On top of that came the traffic those extra atomics generate on a hot per-record
line, and helper work that duplicated the owner's. Removing the helpers removes
the co-installer, and with it the hazard, the latch, the extra compare-and-swaps,
and the traffic. What remains is one driver per
transaction and a single residual wait---the coarse, per-transaction settle-wait
of a foreign owner. That is the sole-driver engine, and it wins the measurement
its predecessor lost.

\subsection*{Contributions}

The paper claims two things, drawn together in \cref{sec:relclaims}: that a
sole-driver facility needs \textbf{no RDCSS}---no restricted double-compare
single-swap, the double-word primitive that normally guards a descriptor
install---and that the facility as a
whole---sole-driven, bounded-blocking, refcount-free to reclaim, usable in
place---is a combination of otherwise-known parts we have not found assembled
before, of which that first result is the linchpin. Four contributions build and
back them:

\begin{itemize}
\item \textbf{The sole-driver \mcas facility} (\cref{sec:facility}): a frozen record
      set, a tri-state status word, and descriptor-naming installs, driven by
      the owner alone. Its main result is a \emph{negative} one---the facility
      needs \textbf{no RDCSS}. A plain
      value-CAS install is A-B-A-safe by construction under one driver, because
      the re-plant hazard it resolves requires a foreign driver of one's own
      records, which the regime excludes (\cref{sec:norddcss}); and its owner-only
      settle makes a terminal transaction a \emph{refcount-free} reclamation point
      (\cref{sec:settle}). The same descriptor carries exclusive and shared
      records together, so exclusion is a property of the \emph{slot} rather than
      of the facility, and the kind never reaches a reader (\cref{sec:kinds}).
      This is the first claim, and three elements of the second.

\item \textbf{The progress argument} (\cref{sec:progress}): deadlock-freedom by
      construction from a sorted install order; \emph{bounded-blocking} defined
      and placed in the progress taxonomy---not obstruction-free, because it
      waits, but starvation-free through a FIFO-fair mutex escalation whose
      budget scales with an operation's cost: the second claim's bounded-blocking
      element.

\item \textbf{A measured design-space map} (\cref{sec:variations}): helping,
      abort-others, age-as-resolution, DCAS (double compare-and-swap)/RDCSS, and
      pure lock-freedom---each built or considered, each rejected, each with the
      throughput or portability reason that decided it. This is not a third claim
      but the \emph{evidence} behind the progress argument above, and---because
      disclosing an implemented-then-rejected approach dates it as prior
      art---the defensive disclosure of the design space itself.

\item \textbf{Logical deletion, applied} (\cref{sec:tombstone},
      \cref{sec:structures}): per-slot atomicity is not operation atomicity, so a
      \emph{usable} concurrent structure needs more than a wide commit. We reuse a
      known technique---a logical-deletion tombstone that update transactions
      treat as freezing a node---and disclose the design rationale that suits it
      to this facility: because the freeze is updater-only coordination, a reader
      gets a frozen node for free without consulting the tombstone, so it stays
      off the reader's path. A single guarded read folds the one remaining
      liveness precondition into the commit. The technique itself is old---one element
      of the combination (\cref{sec:relclaims}), not a novelty of its own---and
      with it, the companion paper's bidirectional list and hash list are shown
      under concurrent writers.
\end{itemize}

\paragraph{Scope.} This paper is about write-set atomicity under concurrency and
what a mutating structure needs to be usable on top of it. The read-side
programming model it forward-references---the read-policy taxonomy, the
general à-la-carte guard facility, and the conflict and aging
\emph{declarations} that tune escalation---is the subject of the next paper in
the series; this one introduces exactly the one guard its own worked example
forces, and no more. The comparative evaluation---tail latency, and
comparison against existence, RLU, and MV-RLU---is later still.
The facility and the companion \fliplatch share the model and the status
vocabulary---and, at the pinned commit, one descriptor and one status
word---so that an embedder that outgrows exclusion changes which discipline a
slot records under, not how it reasons about a commit.

% ===========================================================================
\section{Background: what exclusion was buying}
\label{sec:background}

The companion paper's model is assumed and only its vocabulary is recalled here.
A \emph{record} is a triple \record{s}{o}{n}: slot address $s$, the value $o$ it
must hold at the commit point, and the value $n$ it takes afterwards. A
\pseudotxn is a frozen set of such records over pairwise-distinct slots, made
visible by a single commit step. The write set is atomic---every record resolves
to $o$ before the commit and to $n$ after---and no slot and no instant is ever
torn. A reader is not a transaction of any kind: it meets the \pseudotxn one slot
at a time, so a traversal may straddle a commit, but it observes a monotone
sequence and never new-then-old.

The \fliplatch bought all of this cheaply by requiring writer exclusion. Under
exclusion there is nothing to contend for: a commit cannot fail, there is no
install-time compare-and-swap, no conflict detection, and no abort. The whole
mechanism is one shared selector word that a single release store flips from
$0$ to $1$. Exclusion is what makes that legal, and the paper stated it as a
precondition of the mechanism rather than a caveat on it.

Remove exclusion and three things that exclusion was silently handling come due.

\paragraph{Contended installation.} Two writers may now want to change the same
slot at the same time. The facility must let one of them install its intent on the
slot, let the other discover that intent rather than corrupt it, and resolve
which of them a concurrent reader sees---all without a lock. This is the facility's
core job, and \cref{sec:facility} is how it does it: an installed record names a
descriptor, and every reader and every writer resolves the slot through that
descriptor's status word rather than through the slot's plain value.

\paragraph{Per-slot atomicity is not operation atomicity.} A \kcas makes $k$
named slots change together, but a structural operation's correctness can depend
on a slot it does \emph{not} change. Deleting a node and inserting beside it can
touch \emph{different} words and both succeed, and the result is a lost node or a
use-after-free---not because either commit tore, but because the two operations
were never contending for a common slot in the first place. Making a concurrent
list usable therefore needs a discipline layered on the bare facility, and
\cref{sec:tombstone} supplies it: logical deletion.

\paragraph{Install-time A-B-A.} The third hazard is the subtlest, and it is the
one the sole-driver decision turns on: it does not arise from concurrency as
such, but from a specific way of \emph{answering} the first hazard. If threads
help each other install a transaction's records, then a record can be installed
by a thread that is not its owner, and a stale second install can resurrect a
retired descriptor after the slot's value has cycled back to what it names.
\Cref{sec:norddcss} shows that this hazard is helper-induced, so refusing to help
dissolves it rather than merely defending against it---which is why the facility
needs no RDCSS.

% ===========================================================================
\section{The Sole-Driver MCAS Facility}
\label{sec:facility}

\subsection{Records, status, and resolution}
\label{sec:status}

A transaction is a frozen set of records plus one tri-state \emph{status word}:
%
% One state with two exits, drawn rather than written out.  The list form said
% "UNDECIDED -> SUCCEEDED (commit)  UNDECIDED -> FAILED (abort)", which names the
% start state twice and so reads as two transitions that happen to share a name
% rather than one word leaving one state by one of two doors.  Exactly one of
% these arrows is ever taken.
%
% The terminal boxes are coloured the colour of the value they resolve to --
% green for new, grey for old, matching fig:tristate's targets -- so the two
% pictures agree without the reader being told they do.
\begin{center}
\begin{tikzpicture}[
  font=\small,
  st/.style   = {draw, rounded corners=2pt, minimum height=7mm, inner xsep=5pt},
  arr/.style  = {-latex, thin},
]
\node[st, very thick, fill=orange!18] (u)                {\UNDECIDED};
\node[st, fill=green!12]              (s) at (4.6,  0.8) {\SUCCEEDED};
\node[st, fill=gray!15]               (f) at (4.6, -0.8) {\FAILED};
\draw[arr] (u.east) -- node[above, sloped, font=\scriptsize] {commit} (s.west);
\draw[arr] (u.east) -- node[below, sloped, font=\scriptsize] {abort}  (f.west);
\end{tikzpicture}
\end{center}
The status word is written once, to a terminal value, and never written back.
Installing a record parks the record's own tagged address in its slot---this is
the descriptor-naming step. \emph{How} that park is placed depends on one further
field the companion paper's records did not carry, the record's \emph{kind}: a
slot with concurrent writers is planted with a CAS, a slot the embedder holds a
lock over is parked with a plain store. \Cref{sec:kinds} is that distinction; it
is a writer-side matter only, and nothing below here depends on it. A
reader or a writer that loads a slot holding a parked record reads that record's
status: \SUCCEEDED resolves to the record's new value, \UNDECIDED or \FAILED to
its old value. Resolution goes through the status word, so \textbf{a slot's plain
value is never load-bearing}: one transaction's new value is the next
transaction's old value, and only the descriptor's identity disambiguates them.

\Cref{fig:tristate} is the mechanism in one picture, and it is deliberately the
companion paper's diagram with the selector replaced by a status word.

% The sole-driver descriptor, installed: N slots -> N records -> ONE status word.
%
% DELIBERATELY the same picture as the companion paper's fig:fliplatch, in the
% same visual vocabulary (identical node styles, identical left-to-right layout:
% slots -> indirection -> one shared word -> resolution).  The series claims in
% prose that the two papers realize ONE model with the engine swapped; a reader
% who has seen P1's figure should recognize this one and see that exactly one
% thing changed -- the shared word gained a third state.  Do not "improve" the
% styling independently of P1's figure; the sameness IS the argument.
%
% The one thing this must show that P1's cannot: TWO of the three states resolve
% identically.  Hence the two arrows on the right and not three -- the gray one
% is labelled with both states that take it, so the convergence is visible rather
% than stated.  That asymmetry is the non-obvious part of sec:status.
%
% Layout note: annotations are anchored to nodes, never to guessed absolute
% coordinates.  An earlier draft placed two text blocks by hand and they
% overprinted each other in the rendered PDF.  If you add anything here, render
% the page and LOOK at it before believing the build's silence.
\begin{figure}[t]
\centering
\begin{tikzpicture}[
  font=\small,
  slot/.style   = {draw, minimum width=15mm, minimum height=6mm, fill=gray!8},
  proxy/.style  = {draw, minimum width=26mm, minimum height=9mm, fill=blue!6},
  target/.style = {draw, rounded corners=2pt, minimum width=11mm, minimum height=5.5mm},
  old/.style    = {target, fill=gray!15},
  new/.style    = {target, fill=green!12},
  sel/.style    = {draw, very thick, minimum width=17mm, minimum height=8mm, fill=orange!18},
  arr/.style    = {-latex, thin},
  dep/.style    = {-latex, thin, densely dotted, gray!70!black},
]

% ---- slots (the words the reader loads) -------------------------------------
\node[slot] (s1) at (0, 1.6)  {\code{s$_1$}};
\node[slot] (s2) at (0, 0)    {\code{s$_2$}};
\node[slot] (s3) at (0, -1.6) {\code{s$_3$}};
\node[above=1.5mm of s1, align=center, font=\footnotesize\itshape] {slots\\(tagged)};

% ---- records ---------------------------------------------------------------
\node[proxy] (p1) at (3.9, 1.6)  {$\{o_1,\, n_1\}$};
\node[proxy] (p2) at (3.9, 0)    {$\{o_2,\, n_2\}$};
\node[proxy] (p3) at (3.9, -1.6) {$\{o_3,\, n_3\}$};
\node[above=1.5mm of p1, align=center, font=\footnotesize\itshape] {records};

\draw[arr] (s1) -- (p1);
\draw[arr] (s2) -- (p2);
\draw[arr] (s3) -- (p3);

% ---- the shared status word: the whole point -------------------------------
\node[sel] (st) at (8.3, 0) {\code{status}};
\node[above=1.5mm of st, align=center, font=\footnotesize\itshape] {descriptor\\(one word)};

\draw[dep] (p1.east) -- (st.north west);
\draw[dep] (p2.east) -- (st.west);
\draw[dep] (p3.east) -- (st.south west);

% ---- resolution: two arrows, because two states share a destination --------
% Labels are SLOPED, for the reason P1's figure records: these arrows rise and
% fall about 20 degrees, and a horizontal label placed above or below its own
% midpoint is crossed by the arrow's line.  Verified by rendering the page.
\node[new] (n1) at (12.9,  1.45) {$n_1$};
\node[old] (o1) at (12.9, -1.25) {$o_1$};
\draw[arr, thick] (st.east) --
  node[above, sloped, font=\scriptsize, pos=0.5] {\SUCCEEDED} (n1.west);
\draw[arr, gray!70!black] (st.east) --
  node[below, sloped, font=\scriptsize, pos=0.5] {\UNDECIDED\ or \FAILED} (o1.west);

% ---- the terminal write ----------------------------------------------------
% Drop is 0.95, not 0.55: the sloped UNDECIDED/FAILED label descends into this
% column, and at the shorter drop its tail and this block's top line came within
% a hair of touching.  Rendered and checked.
\draw[-latex, very thick, orange!70!black] (st.south) -- ++(0,-0.95)
  node[below, align=center, font=\footnotesize, text=black]
  {\textbf{written once, to a terminal value}\\
   \UNDECIDED $\to$ \SUCCEEDED\ (commit)\\
   \UNDECIDED $\to$ \FAILED\ (abort)};

\end{tikzpicture}
\caption{The sole-driver descriptor, installed---the companion paper's figure
with its selector replaced by a status word. Each slot holds a \emph{tagged}
pointer to a record rather than its target; each record carries
$\{\mathit{old},\mathit{new}\}$ and reads the one shared status word, so a single
store decides the whole set. What the third state buys is on the right:
\SUCCEEDED\ resolves every record to its new target, while \UNDECIDED\ and
\FAILED\ both resolve to old. An installed-but-undecided set and an aborted set
are therefore indistinguishable to a reader, and an abort leaves the structure
exactly as it was found. As in the single-writer case there is no instant at
which the set is part-old and part-new.}
\label{fig:tristate}
\end{figure}


This is the shape the companion paper inherited from practical \mcas
\citep{harris2002mcas, fraser2004lockfreedom}: a descriptor of records, a status
word whose transition decides the outcome, and resolution routed through the
status rather than through the slot. What is new here is not the shape but what
becomes unnecessary once the single-driver discipline is added to it.

\subsection{Sole-driver: no helping, no stealing, not even reads}
\label{sec:soledriver}

\begin{definition}[Sole-driver discipline]
\label{def:soledriver}
A transaction is installed, decided, and settled by its \emph{owner} alone. No
foreign thread ever plants, drives, evicts, or steals another transaction's
records. A thread that meets a foreign in-flight record resolves it to the
transaction's logical value, or---when it needs the slot itself settled before it
can act on it---waits a bounded spin for the owner to settle it; either way it
proceeds without ever driving the record.
\end{definition}

The discipline is total, and its most easily overlooked edge is the read side. A
reader that lands on an \UNDECIDED foreign record does not help decide it: it
resolves the record to its \emph{logical} value---the old, the value the
transaction takes if it aborts---at once, without waiting on it, and moves on. The
facility's own read accessor states it as an imperative: \emph{never drive} the
foreign transaction. This closes an objection a careful reader will otherwise
raise---``but surely your read path helps''---and it is worth stating outright
because it is what makes the next subsection's argument airtight: there is no
foreign \emph{installer} of a transaction's records, and no foreign \emph{driver}
of any kind, reads included.

A bounded wait is a bias, not a guarantee. It buys the common case, in which the
owner settles promptly; it does not close the window against an owner that has
been preempted mid-transaction. What closes that window is not the wait but the
escalation of \cref{sec:progress}. Keeping the wait and the progress guarantee
separate is deliberate: the wait is an optimization, and the correctness of the
facility does not rest on it.

\subsection{No double-word install}
\label{sec:norddcss}

A value-CAS install is the natural way to place a record: compare-and-swap the
slot from the old value the caller saw to the record's tagged address. The
classical objection is an A-B-A: between the caller's read and its install, the
slot could move to some other value and back to the old one, and a naive install
would then succeed against a slot whose meaning had changed underneath it. The
standard defence is RDCSS---a restricted double-compare-single-swap that
conditions the install on a second word, typically the descriptor's own
status---and it is exactly what Harris, Fraser and Pratt introduced to make
practical \mcas safe~\citep{harris2002mcas}.

It is worth being concrete about that hazard, because what follows is a claim
that it \emph{cannot arise}, not that it is improbable. Let a transaction $T$
hold a record $r$ on slot $s$, and let $H$ be some thread helping $T$ install it:

\begin{enumerate}\setlength{\itemsep}{1pt}
\item $H$ reads $s$ and finds the value $A$ that $r$ expects.
\item $H$ stalls---descheduled, preempted, anything---before its
      compare-and-swap.
\item Meanwhile $T$ is driven to completion by its owner or by other helpers. It
      decides, settles every slot back to a plain value, and is retired. Nothing
      names its descriptor any longer, so a grace period later that memory is
      handed back.
\item Unrelated work changes $s$ and changes it back, so $s$ holds $A$ once more.
      A next pointer cycling $A \rightarrow X \rightarrow A$ does this, and so
      does an embedder recycling a value.
\item $H$ resumes and performs its compare-and-swap. \emph{It succeeds}, because
      $s$ does hold $A$.
\end{enumerate}

\noindent Slot $s$ now names a record belonging to a transaction that finished
long ago. A reader resolving $s$ follows $r$ to $T$'s status word, finds it
terminal, and takes $T$'s value---from a transaction already settled, and whose
descriptor may since have been freed. The slot has been re-parked by a thread
that had no business writing it, which is why the failure is not merely a stale
read but a use-after-free waiting to happen. RDCSS closes step~5 by conditioning
the install on a second word: the descriptor's status must still be \UNDECIDED,
and $T$'s no longer is.

Now notice which step carries the weight. Steps~1, 4 and~5 are ordinary---a read,
a value recurrence, a compare-and-swap that finds what it expected. Step~3 is
just a transaction completing. Everything hinges on step~2: a thread that had
resolved to install one of $T$'s records and had not yet done so when $T$
finished. That thread cannot be $T$'s owner. An owner plants its whole prefix and
only then decides, so at the instant $T$ becomes terminal its owner has no plant
outstanding---there is nothing left for it to stall before. The stalled installer
of step~2 is therefore a \emph{second} driver of $T$, and a second driver is what
helping introduces and what \cref{def:soledriver} forbids. The property states
that.

The sole-driver facility omits RDCSS. It installs with a bare value-CAS, and that
is sufficient. The reason is precise, and it is the paper's strongest claim.

\begin{property}[Install-A-B-A is helper-induced]
\label{prop:abainstall}
The install-time A-B-A that RDCSS prevents requires a \emph{second, stalled
driver of the same transaction}: a thread that installs one of a transaction's
records late, after the slot's value has recurred, resurrecting a descriptor its
owner has already retired. Under \cref{def:soledriver} no such driver exists.
The owner plants its whole prefix before it decides, plants each record exactly
once while the transaction is \UNDECIDED, and no foreign thread ever re-plants,
drives, or settles that record. So a stale second plant cannot occur, and any
recurrence of a slot's plain \emph{value}---a next-pointer cycling
$B \rightarrow X \rightarrow B$ with $B$ a live, never-freed successor, a counter
revisiting a number, an embedder's own reuse---is inert. A losing install is a
plain CAS failure the caller handles by re-reading; it never resurrects anything.
\end{property}

This is applied ownership logic, not a new primitive. The safety of a
disjoint-access-parallel operation that only one thread may drive is the
classical observation of Israeli and Rappoport~\citep{israeli1994disjoint}; the
contribution here is the recognition that a \pseudotxn's install inherits exactly
that safety once helping is removed. The linchpin is that RDCSS's purpose is tied
to helping in the first place: a recent analysis of \kcas frames RDCSS precisely
as the machinery that keeps an install safe \emph{when the operation may be
completed by a helping thread}~\citep{guerraoui2020mcas}. A helping \kcas states
the same dependency from the inside: PathCAS adopts the double-word install
specifically because it ``avoids ABA problems that could otherwise be introduced
by this lock-free helping,'' preventing a helper from installing into a
descriptor whose outcome is already decided and so \emph{resurrecting a completed
operation}~\citep{brown2022pathcas}. That is \cref{prop:abainstall} read in the
other direction, by authors who kept the helpers. Remove the helping and
the machinery it justified goes with it. We are not aware of prior work that
states the elimination this way---that the sole-driver regime makes RDCSS
unnecessary rather than merely cheaper.

The facility does still require the ordinary things a value-based \mcas requires,
and none of them concerns slot values: the embedder's per-record tag kept free
in every legitimate slot value (\cref{sec:constraints}), pairwise-distinct slots
per transaction, and---the reclamation obligation---that a descriptor
reachable through a slot is freed only after a grace period, because a reader may
still dereference it. This is value-CAS atomicity: a record is validated to hold
its old at the linearization point, not to have been stable throughout. An
embedder that needs stability across an interval layers its own versioning on
top, as with any value-based \mcas. Memory safety holds unconditionally.

% -----------------------------------------------------------------------
% NOTE (drafting; retained deliberately -- see README framing rule).
% This is the ENGINE-CLAIMS LEDGER item 2, the centerpiece. Keep it framed as
% (a) a NEGATIVE result -- we removed a mechanism -- and (b) applied ownership
% logic, NOT a new primitive. Israeli-Rappoport for the ownership base;
% Guerraoui 2020 as the linchpin that ties RDCSS to helping. "Apparently
% unpublished" is the register: "we are not aware of", never "novel"/"first".
% Do NOT let an editing pass strengthen the last sentence into a bare claim.
% -----------------------------------------------------------------------

\subsection{Owner-only settle, and where reclamation becomes safe}
\label{sec:settle}

After a transaction is decided, its owner \emph{settles}: it rewrites each record
it planted back to a plain value---the new value on \SUCCEEDED, the old on
\FAILED---so the indirection is gone and later readers resolve directly. Settle
is owner-only and never waits. Because nothing foreign ever planted into those
records, and the owner converts exactly the prefix it planted, one property holds
the instant settle returns: \textbf{no proxy of the transaction names any slot}.

That property is the whole of the reclamation argument, and it is important to
claim exactly it and no more. Reclaiming descriptors after a grace period is not
new---it is the standard \rcu treatment of any reader-reachable object, and it
has been named and benchmarked as an \mcas reclamation baseline over this very
library~\citep{arbelraviv2017descriptors}. What the sole-driver regime adds is
narrower and is a \emph{reachability} statement: because no foreign thread can
plant, drive, or steal a transaction's records, a terminal-and-settled
transaction is a point at which nothing can any longer make its descriptor
reachable, so the descriptor may be handed to \code{call\_rcu()} with no
reference count and no further interlock. The safety of the free is the ordinary
\rcu grace period; what is owner-derived is only the certainty that the free
point has been reached.

% -----------------------------------------------------------------------
% NOTE (drafting; retained deliberately -- see README framing rule).
% HIGH overclaim risk (ENGINE-CLAIMS LEDGER item 5). Do NOT claim "reclaim
% descriptors via RCU" -- that is Arbel-Raviv--Brown's named, benchmarked
% baseline over liburcu. Claim ONLY the refcount-free REACHABILITY point:
% terminal-and-settled means nothing foreign can resurrect the records, so it
% is a safe call_rcu point WITHOUT a refcount. The grace period does the
% freeing; sole-driver only certifies the point. Keep this distinction sharp.
% -----------------------------------------------------------------------

\subsection{Two install disciplines, one linearization point}
\label{sec:kinds}

% Two install disciplines in, one resolve path out.
%
% This figure exists to DRAW prop:kindinvisible rather than argue it: every
% difference between the kinds is on the ARROW from slot to record, and nothing
% past the record can tell which one parked a slot.  If a reader takes only that
% away, the figure has done its job.
%
% Same visual vocabulary as fig:tristate and P1's fig:fliplatch -- deliberately.
%
% The slots are drawn IDENTICALLY on purpose.  It is tempting to colour the
% exclusive one differently, but that would be untrue to the mechanism: a slot's
% kind is not a property the word carries, it is a property of the WRITER's
% discipline, so it belongs on the arrow and nowhere else.
%
% Deliberately NOT drawn: the install ORDER (shared records first, then the
% exclusive parks, then the status store, then settle).  That is sequencing, it
% belongs in prose, and crowding it in here would bury the single point above.
%
% Layout note: an earlier draft ran a rotated "the kinds differ here" label down
% the middle and it overprinted the arrow labels.  The two shared slots sit close
% together and the exclusive one sits apart, which separates the kinds by
% POSITION instead, leaving the top clear for the failure edge.  Render the page
% and look at it before adding anything.
\begin{figure}[t]
\centering
\begin{tikzpicture}[
  font=\small,
  slot/.style   = {draw, minimum width=15mm, minimum height=6mm, fill=gray!8},
  proxy/.style  = {draw, minimum width=26mm, minimum height=9mm, fill=blue!6},
  target/.style = {draw, rounded corners=2pt, minimum width=11mm, minimum height=5.5mm},
  old/.style    = {target, fill=gray!15},
  sel/.style    = {draw, very thick, minimum width=17mm, minimum height=8mm, fill=orange!18},
  arr/.style    = {-latex, thin},
  dep/.style    = {-latex, thin, densely dotted, gray!70!black},
  fail/.style   = {-latex, thin, densely dashed, red!60!black},
]

% ---- slots: drawn identically; the kind is not visible in the slot ----------
\node[slot] (s1) at (0,  1.7) {\code{s$_1$}};
\node[slot] (s2) at (0,  0.5) {\code{s$_2$}};
\node[slot] (s3) at (0, -1.5) {\code{s$_3$}};
\node[left=1.5mm of s1, font=\scriptsize\itshape] {shared};
\node[left=1.5mm of s2, font=\scriptsize\itshape] {shared};
\node[left=1.5mm of s3, font=\scriptsize\itshape, align=right] {exclusive\\(lock held)};

% ---- records ---------------------------------------------------------------
\node[proxy] (p1) at (4.4,  1.7) {$\{o_1,\, n_1\}$};
\node[proxy] (p2) at (4.4,  0.5) {$\{o_2,\, n_2\}$};
\node[proxy] (p3) at (4.4, -1.5) {$\{o_3,\, n_3\}$};

% ---- the two installs: the ONLY place the kinds differ ---------------------
\draw[arr] (s1) -- node[above, font=\scriptsize] {CAS} (p1);
\draw[arr] (s2) -- node[above, font=\scriptsize] {CAS} (p2);
\draw[arr, thick] (s3) --
  node[above, font=\scriptsize] {plain store}
  node[below, font=\scriptsize] {cannot fail} (p3);

% The failure edge belongs to the CAS installs and to nothing else.  Anchored to
% the arrow's own midpoint, drawn up into space kept clear for it.
\draw[fail] (2.2, 1.95) -- ++(0.5, 0.75)
  node[right, font=\scriptsize, text=red!60!black] {a peer changed it: \emph{abort}};

% ---- one status word -------------------------------------------------------
\node[sel] (st) at (8.8, 0.2) {\code{status}};
\draw[dep] (p1.east) -- (st.north west);
\draw[dep] (p2.east) -- (st.west);
\draw[dep] (p3.east) -- (st.south west);

% ---- one resolve path, identical for both kinds ----------------------------
\node[old] (res) at (12.6, 0.2) {$o_i$ / $n_i$};
\draw[arr, thick] (st.east) -- node[above, font=\scriptsize] {resolve} (res.west);

% Short, and tucked directly under the resolve arrow it labels.  Two earlier
% placements collided with the dotted arrows converging on `status` from the
% records -- that whole wedge of space is crossed by them and is not usable for
% text.  Under the resolve arrow, right of `status`, nothing crosses.  The long
% form of this sentence lives in the caption, where it cannot collide with
% anything.
\node[font=\footnotesize\itshape, text=gray!45!black]
  at (10.85, -0.5) {identical for both kinds};

\end{tikzpicture}
\caption{Two install disciplines, one resolve path. A slot with concurrent
writers is planted with a compare-and-swap, which a peer's intervening change
makes fail---the only source of a contention abort. A slot the embedder holds a
lock over is parked with a plain release store, which cannot fail. The slots are
drawn identically because they \emph{are} identical: the kind is a property of
the writer's discipline, not something the word carries, so it lives on the
arrow and nowhere else. Past the records the two are one object, and a reader
resolving a slot cannot determine which discipline parked it. A transaction with
no shared record therefore has nothing in it that can abort.}
\label{fig:kinds}
\end{figure}


Everything above assumed every slot has concurrent writers. Real structures are
rarely so uniform: an edit often spans slots the embedder already holds a lock
over and slots it does not, and until recently that edit had no single commit
spanning both---it had to bridge two facilities with an external lock. The
facility removes the bridge by making the install discipline a property of the
\emph{record}, chosen per slot:

\begin{description}\setlength{\itemsep}{2pt}
\item[Shared (\mw).] The slot has concurrent writers. The install is the
      sole-driver CAS of \cref{sec:norddcss}: a writer that changed the slot
      fails the CAS, and the transaction aborts.
\item[Exclusive (\sw).] The embedder holds a lock over the slot, so no peer can
      race it. The install is a plain release-store park. It cannot fail.
\end{description}

The two differ \emph{only} in how a record is placed and settled. They share the
resolve header---the old value, the new value, and the pointer to the one status
word---so a slot resolves identically whichever kind parked it, and the single
release store of that status word is the linearization point for all of them at
once. Three consequences follow, and they are the whole of what the distinction
buys.

\begin{property}[The kind is writer-side only]
\label{prop:kindinvisible}
A reader loads a slot, finds a tag, and resolves through the status word exactly
as in \cref{sec:status}. Nothing in that path consults the kind, and no reader
can determine which discipline parked a slot. Generality on the writer side is
therefore free on the read side---the property the whole series is organized
around.
\end{property}

That this is even available is worth stating, because the obvious layout does not
give it. If the exclusive discipline had kept its own proxy layout and the shared
one its descriptor, a single tag bit would name two incompatible objects and a
reader would have to know which resolver to run---a genuine obstacle, and the
reason the two lived in separate facilities. Sharing the resolve header dissolves
it: one layout, one resolver, and the incompatibility turns out to have been an
artefact of representation rather than of the problem.

\emph{Abort is a shared-slot phenomenon.} The commit installs the shared records
first and the exclusive parks after, so an abort happens before any park is
placed and no park is ever rolled back---only the planted prefix is settled back.
A write set with no shared record therefore \textbf{cannot contention-abort at
all}: it is the companion paper's \fliplatch commit, recovered exactly, as the
degenerate case of this one. Progress is thus a property of the write set rather
than of the facility, which \cref{sec:boundedblocking} takes up.

\emph{Mixing is paid for only when it happens.} Grouping the shared records ahead
of the exclusive ones is the one cost a uniform write set would not otherwise
bear, so the commit tracks how many records are of each kind and does the
regrouping only when a write set genuinely holds both. An all-shared commit costs
what it did before the kinds existed; an all-exclusive one skips the multi-writer
path entirely, through a second commit entry point an embedder uses when it knows
its write set carries no shared record.

\subsection{Constraints}
\label{sec:constraints}

An installed record is recognized by a \emph{tag}, and---as in the companion
paper---the tag is chosen \emph{per record} by the embedder rather than fixed by
the facility: bit~0 in the structures shown here, a reserved nibble or high bits
where a slot's own encoding has already spent its low bit. Whatever the embedder
picks, every value that slot may legitimately hold must keep the tag pattern
clear. The facility's own fixed requirement is narrower and always met: a record's
address must carry the tag pattern in free low bits, which its alignment
guarantees.
The record set is frozen at commit---membership is fixed before installation
begins---and the slots of one transaction are pairwise distinct. These contracts
are discharged by the embedder's encoding, by construction or not at all.

The kinds of \cref{sec:kinds} add one more, and it is the sharpest of them
because it cannot be checked. A slot is exclusive \emph{or} shared, and which one
must hold \textbf{globally}: every transaction that records that slot, on every
thread, must agree. The asymmetry is what makes this dangerous. Naming a slot
shared is always safe---a CAS install is correct whether or not anyone else
touches it---whereas naming it exclusive is a \emph{promise}, and the plain park
is correct only while the embedder's lock really does exclude every other writer.
A transaction that treats a slot as exclusive while another treats it as shared
lets the second's CAS race the first's unguarded park, which tears or loses a
write. The facility cannot detect this, because each transaction sees only its
own records; it is the embedder's invariant, exactly as ``the same datum is
guarded by the same lock everywhere'' is. A quiescent transition---a slot shared
for one phase and exclusive for another, with a grace period between---is
legitimate; concurrent disagreement is not. When in doubt, name it shared. Within
a single transaction the contradiction is at least visible, and the commit
resolves it in the safe direction, promoting the slot to the CAS install rather
than trusting a promise the record set itself contradicts.

% ===========================================================================
\section{Progress and the Coarseness Thesis}
\label{sec:progress}

\subsection{Deadlock-freedom by construction}
\label{sec:deadlock}

A committer that waits on a foreign owner while holding installs of its own could
deadlock if two committers waited on each other. The facility prevents it not with
detection but with a total order: records install in one global order, sorted by
slot address. A committer therefore holds only \emph{lower} slots while it waits
on a shared one, and the owner it waits on---having itself reached that shared
slot---already passed every lower slot without blocking, so it holds none the
waiter holds. There is no cycle to close, so there is no deadlock.

The sort is paid only when it is needed. A first, optimistic attempt installs
flat, in caller order, and never waits: the moment it meets a foreign record it
bails and fails fast, which cannot deadlock because it never holds-and-waits. Only
when an attempt escalates to the blocking, waiting path does it first sort its
records by address, buying the ordering the wait needs. Age~0 is unsorted and
non-blocking; age~1 and above are sorted and blocking.

The order covers the shared records alone. An exclusive park waits on nobody---by
the promise that makes it exclusive there is no foreign owner to wait on---so it
can close no cycle, and ordering it would buy nothing.

\subsection{Bounded-blocking, defined}
\label{sec:boundedblocking}

\begin{definition}[Bounded-blocking]
\label{def:boundedblocking}
An operation is \emph{bounded-blocking} if the only way it waits on another
thread is a wait of bounded duration on that thread to complete a bounded amount
of work, after which the waiter either proceeds or escalates to a mechanism that
guarantees its own progress.
\end{definition}

``Bounded-blocking'' is our descriptive term, not an established progress class,
and it is worth being exact about where it sits. The facility \emph{waits} on a
foreign owner, so it is blocking, and therefore it is not
obstruction-free~\citep{herlihy2003obstruction}: an obstruction-free operation
completes if it runs in isolation for long enough, and this one can be made to
wait even then. That is a real concession and we make it plainly. What the facility
does provide is stronger than deadlock-freedom: through the escalation of
\cref{sec:escalation}, a starved committer eventually acquires a FIFO-fair lane
and completes, so under a fair scheduler the facility is
\emph{starvation-free}~\citep{herlihy2011progress}. Bounded-blocking is thus the
conjunction of a bounded common-case wait and a starvation-free fallback, and it
sits below lock-freedom and above deadlock-freedom in the taxonomy.

One refinement follows from \cref{sec:kinds}, and it is worth stating because it
narrows the concession rather than restating it. What waits, and what can abort,
is a \emph{shared} record; an exclusive one does neither. So bounded-blocking
classifies a \emph{write set}, not the facility: a write set carrying no shared
record commits without waiting on anyone and without any contention abort, which
is the companion paper's progress, unchanged. A transaction pays the
blocking-progress price on exactly the slots that have a concurrent writer, and
an embedder that can exclude writers from some of its slots pays it on fewer of
them. This is the coarseness thesis at a finer grain than a whole facility: the
synchronization is as coarse as the sharing, and no coarser.

\subsection{Adaptive coarseness: the fair-mutex escalation}
\label{sec:escalation}

The optimistic path is bounded-blocking but not by itself starvation-free: a
large transaction can be defeated indefinitely by a stream of small ones that
keep changing a slot between its read and its install. The facility handles this by
escalating a losing handle into a per-domain \emph{fair mutex}---an MCS-style
queue lock~\citep{mellorcrummey1991mcs}---and publishing a flag so that future
transactions in the domain funnel through the same lane. Inside the lane the
escalated transaction contends only with the finite in-flight set, bounded by the
thread count, and commits within a bounded number of retries while holding its
turn. Progress is guaranteed with no quiescence: no \code{synchronize\_rcu()},
no drain.

The escalation is reactive, and its threshold is the interesting part. A handle
takes the lane not because its transaction is large but because it is
\emph{losing}: it has retried past a budget that its own cost earns it. The
budget scales with cost---measured as the transaction's read count plus its
write-set size, everything the attempt touched rather than only the edges it
commits---for a
reason that is itself a small throughput argument. Escalating funnels every
writer in the domain through one serial lane, so the cost of escalating is the
length of the critical section it serializes. A single-edge operation commits
with a bare CAS: serializing it is nearly free, the lane is just orderly queueing,
and escalating early turns wasted CAS collisions into turns. A traversal mutator
is the opposite: funnelling those destroys the $N$-way parallelism of the
optimistic path, which aborts far more often but retries in parallel.

How much that costs is not one number. The lane is a single global critical
section, so its throughput is flat in the writer count while the optimistic path
scales with it. Their ratio therefore grows with the writer count, and any
single factor is only a reading taken at one point on that curve. Measured on a
three-structure skip-list move, optimistic against a build that escalates on its
first attempt:

\begin{center}
\footnotesize
\setlength{\tabcolsep}{5pt}
\begin{tabular}{@{}l rrrrrrrrrr@{}}
\toprule
writers & 1 & 2 & 4 & 8 & 16 & 32 & 48 & 80 & 128 & 192 \\
\midrule
optimistic $/$ lane & $1.0\times$ & $2.3\times$ & $26\times$ & $49\times$ &
  $95\times$ & $180\times$ & $224\times$ & $970\times$ & $941\times$ &
  $1042\times$ \\
\bottomrule
\end{tabular}
\end{center}

\noindent At a single writer the two are identical---an uncontended lane costs
nothing, which is also the check that the always-escalate build is measuring the
funnel rather than itself---and past 80 writers the optimistic path is three
orders of magnitude ahead, plateauing there as the lane's own throughput goes
flat.\footnote{\code{bench\_txn\_3skiplist}, width~3, 3840
keys per skip list, median of nine 2\,s runs with the address-space layout
pinned, on the machine of \cref{sec:varhelping}. The always-escalate arm forces
the lane from the first attempt (a zero budget); it is the limit case, not the
shipped policy, which escalates only a losing handle. Sixty-four writers is
omitted deliberately, and is the one point that is not a number: the serialized
lane has two stable regimes either side of it---$0.161$ million key-moves per
second at 48 writers and $0.046$ at 80, each reproducing to within
$0.1\%$---and at 64 it settles into one or the other from run to run, so its
ratio ranges over roughly $270$--$810\times$. Every other point here reproduces
within $3\%$. The comprehensive study is deferred to the evaluation paper.} So the more a
transaction costs, the longer it should stay optimistic---the budget encodes
exactly that.

\subsection{Preemption: a design lever, not a claim}
\label{sec:preemption}

A lock-free operation \emph{tolerates} preemption at a standing cost: it is
designed so that a thread stopped at any point obstructs no one. A
bounded-blocking operation instead has a preempted owner as its worst case, the
one the escalation exists to rescue. There is a third option that neither
tolerates nor merely rescues: \emph{avoid} the preemption. The Linux kernel has
gained an \code{rseq} time-slice extension that lets a userspace thread defer its
own preemption for a bounded window~\citep{gleixner2025rseqtse}; applied to the
install-through-settle section of a sole-driver commit, it would let an owner
reach settle before it is taken off-CPU, turning the preempted-owner tail from
\emph{bounded} into \emph{rare}.

We are deliberately careful about what is and is not claimed here.

% ---------------------------------------------------------------------------
% NOTE (drafting; retained deliberately -- see README framing rule + SCOPE).
% RSEQ-TSE is NOT in the pinned tree -- verified against rcu-txn.h at the pin:
% no rseq, no time-slice, not even a TODO. So this subsection describes a design
% lever, NOT a measured feature, and the current throughput numbers are the
% PESSIMISTIC case without it. Mathieu's claim boundary (2026-07-18):
%   - claim NOTHING on the mechanism (the mainline rseq-TSE series, Gleixner
%     2025, which the author took part in);
%   - claim NOTHING on the use-pattern "defer preemption of a critical-section
%     holder" -- textbook (schedctl / temporary non-preemption / scheduler-
%     conscious sync); cite the lineage, do not claim it;
%   - DISCLOSE, as defensive prior art, the userspace use for THIS purpose
%     (deferring preemption across the sole-driver install->settle window), and
%     claim only the facility-level OBSERVATION that it thins the tail from
%     bounded to rare -- which is what makes bounded-blocking practically
%     competitive.
% RESOLVED (Mathieu, 2026-07-19): KEEP the subsection (option a) -- explicit
% future/defensive-disclosure, clearly marked unimplemented. The schedctl/TNP
% lineage cites are added and DBLP/patent-verified: schedctl = US Patent
% 5,937,187 (kosche1999preemption) + its unmerged Linux restatement, Aziz LKML
% 2014 (aziz2014preempt); temporary non-preemption = Marsh et al. SOSP
% 1991 (marsh1991firstclass); scheduler activations = Anderson et al. SOSP 1991
% (anderson1991scheduler); lock-holder-preemption avoidance = Dice-Harris
% TRANSACT 2016 (dice2016lockholder).
% ---------------------------------------------------------------------------
The \emph{mechanism} is not ours: it is the mainline kernel series, and the
general pattern of deferring a critical-section holder's preemption so waiters do
not stall is textbook---Solaris's \code{schedctl} preemption
control~\citep{kosche1999preemption} and its later Linux
restatement~\citep{aziz2014preempt}, the temporary non-preemption of
\citet{marsh1991firstclass}, scheduler activations~\citep{anderson1991scheduler},
and lock-holder-preemption avoidance~\citep{dice2016lockholder}. What
we disclose, as prior art rather than claim as invention, is the specific
userspace use---deferring preemption across a bounded-blocking \mcas's publish
window---together with the facility-level observation that it thins the tail. And
we are explicit that this lever is \textbf{not present in the engine measured
here}: the throughput this paper reports is the pessimistic case, \emph{without}
it, so the lever is unclaimed upside rather than a result. Correctness does not
depend on it: it is best-effort, bounded, and not available on every kernel, and
the aging-to-fair-mutex escalation carries the residual either way.

\subsection{The trade, stated plainly}
\label{sec:trade}

The facility gives up lock-freedom's standing guarantee---that no thread's delay
can stall another---and gets, in exchange, a mechanism with fewer synchronization
points, no per-record install state, less hot-cache-line traffic, and a commit
that is a bare CAS per slot in the common case. The concession is a tail: a
preempted owner blocks its waiters until the escalation resolves it. The claim of
this paper is that for the read-dominated structures \rcu serves, and at the
scale where these facilities are used, that is the right end of the trade---and the
next section is the evidence, because the losing end was built too.

% ===========================================================================
\section{Variations and Alternatives}
\label{sec:variations}

% ---------------------------------------------------------------------------
% THIS SECTION IS THE DEFENSIVE PAYLOAD and P2's strongest asset: a MEASURED
% design-space map. Each alternative was built (or seriously considered) and
% then rejected for a throughput or portability reason. Disclosing an
% implemented-then-rejected approach enables it as prior art for free.
%
% MEASUREMENT (SCOPE DECISION #1) -- RESOLVED 2026-07-19, RE-MEASURED 2026-07-28:
% the decisive number is the help+steal -> sole-driver A/B, now +22.0% at 192
% writers (to roughly +30% under contention), from
% efficios-trie-benchmark scripts/p2_engine_remeasure.csv; both arms built from
% engine 97443472 and differing in one flag. Retire commit a07b67ba.
% abort-others (= stealing) is folded into that bundle; age-as-resolution has no
% isolated run and is rejected on the correctness trap. No \UNVERIFIED remain.
% See the long note in sec:varhelping for why the A/B is at 97443472 and for the
% three preparation traps -- read it before refreshing these numbers again.
% ---------------------------------------------------------------------------

The mechanism of \cref{sec:facility} is one point in a design space, and unusually
for such a section the neighbouring points are not hypothetical: several were
implemented and measured, then removed. That is what makes the map defensively
load-bearing---an implemented-then-rejected approach is disclosed and dated
here---and it is what makes the coarseness thesis evidence rather than assertion.
It begins, though, with the premise rather than the alternatives: whether
dropping the exclusion buys anything at all.

\subsection{What dropping the exclusion buys}
\label{sec:writerscaling}

The companion paper's \fliplatch requires writer exclusion. That is not a
property of its implementation that a later version might relax---it is the
precondition the whole mechanism rests on, which is why a commit there can be a
single store with no compare-and-swap, no conflict detection and no abort. So the
question this paper exists to answer has a measurable form: with a second writer,
and a hundredth, what does the difference amount to?

\Cref{fig:writerscale} is that comparison on the bidirectional list, which both
papers use. The single-writer facility is measured with its writers serialized
through a mutex, because that is what its precondition means once a second writer
exists; the alternative is not to run the workload at all. At one writer the two
are within $3\%$, so the transacted facility is not buying its scaling by
degrading the uncontended case---it starts level. From there they separate: the
transacted list reaches $321$ million updates per second at 192 writers, a factor
of $50$ over its own one-writer figure, while the excluded one peaks at its
\emph{first} writer and declines thereafter, ending at a sixth of its own
one-writer throughput on 192 times the hardware. Serialization is not slow growth; it is negative growth, and
a plain mutex list behaves the same way.

Two things about that figure should be read narrowly, and are drawn so they
cannot be read otherwise. The writers are \emph{disjoint by construction}---each
owns its own churn slots, and no two ever touch the same one---so this measures
what exclusion costs when there is nothing to serialize, which is the strongest
form of the objection to it and not a favourable case for us. And there are no
readers at all: the workload is pure mutation, which for a facility built to
serve read-mostly structures is close to an anti-workload. What the figure
supports is exactly one claim, that concurrent writers proceed concurrently. It
is not a comparative evaluation, and the sublinearity it also shows---$26\%$ of
ideal at 192 writers---belongs to this facility as much as the scaling does.

% Writer scaling on the bidirectional list: what dropping the exclusion buys.
%
% This is the paper's PREMISE measured.  Everything else here argues about how
% the facility should be built; this is the evidence that building it at all was
% worth doing, and until now the paper asserted it.
%
% The structure is the companion paper's, deliberately -- the series keeps the
% same worked structures so what changes between papers is the facility and not
% the example.  The single-writer arm is that paper's flip-latch with its writers
% serialized through a mutex, which is not a handicap imposed on it: it is what
% exclusion MEANS once there is more than one writer.
%
% LOG-LOG, and for a reason.  The two curves end up three orders of magnitude
% apart, so a linear y would erase the single-writer curve entirely -- it would
% sit on the axis and the reader would have to take its shape on trust.  The
% ideal-scaling guide is drawn for the same reason it would be tempting to omit:
% on log-log, perfect scaling is a straight line of slope 1, so the gap between
% the measured curve and that line IS the sublinearity, stated rather than
% hidden.  At 192 writers the facility reaches 26% of ideal.
%
% A plain pthread_mutex list is NOT plotted, though it was measured and the data
% file carries it.  Its absolute level is not comparable: the lock engines
% recycle a permanent node in place while the RCU engines defer frees through
% call_rcu, so it starts far above both curves at one writer and would invite
% exactly the wrong reading.  Its SHAPE is the useful part -- flat from two
% writers on, like the flip-latch -- and the caption says so in words instead.
\par\medskip\noindent\begin{minipage}{\linewidth}
\centering
\begin{tikzpicture}
\begin{axis}[
  width=0.92\linewidth, height=5.6cm,
  xmode=log, ymode=log, log basis x=2,
  xlabel={Concurrent writers}, ylabel={Millions of updates/s},
  xmin=0.85, xmax=230, ymin=0.7, ymax=1600,
  xtick={1,2,4,8,16,32,64,128,192},
  xticklabels={1,2,4,8,16,32,64,128,192},
  ytick={1,10,100,1000},
  yticklabels={1,10,100,1000},
  legend pos=north west, legend cell align=left,
  legend style={draw=none, fill=none, font=\footnotesize},
  tick label style={font=\footnotesize},
  label style={font=\footnotesize},
  grid=major, grid style={gray!20},
  axis lines=left,
]
\addplot[no marks, thin, gray!55, densely dotted, domain=1:192, samples=2]
  {6.40*x};
\addlegendentry{ideal scaling}
\addplot[mark=*, mark size=1.5pt, thick, color=blue!65!black]
  table[x=writers, y=mw, col sep=comma] {data/fig-writerscale.csv};
\addlegendentry{\code{rcu-txn-list} (this paper)}
\addplot[mark=square, mark size=1.5pt, thick, color=black!55, dashed]
  table[x=writers, y=sw, col sep=comma] {data/fig-writerscale.csv};
\addlegendentry{\fliplatch $+$ writer exclusion}
\end{axis}
\end{tikzpicture}
\captionof{figure}{Writer scaling on the bidirectional list of
\cref{sec:structures}, the structure the companion paper also uses. The
single-writer \fliplatch is shown with its writers serialized through a mutex,
which is what its exclusion precondition amounts to once a second writer exists.
The two are within $3\%$ of each other at one writer and diverge from there: the
transacted list reaches $321$ million updates per second at 192 writers, a factor
of $50$, while exclusion turns the second writer into a loss and every writer
after it into a smaller one. A plain \code{pthread\_mutex} list, measured but not
drawn, has the same flat shape---serialization is flat whatever sits behind
it---but its absolute level is not comparable, because the lock engines recycle a
node in place where the \rcu ones defer every free through a grace period. The
dotted guide is ideal scaling from the one-writer point: the measured curve
reaches $26\%$ of it at 192 writers, and the distance between them is this
facility's own sublinearity, which no amount of dropping exclusion removes.}
\label{fig:writerscale}
\end{minipage}\par\medskip


\subsection{Helping, and the latch it forced}
\label{sec:varhelping}

The engine once helped. Threads that met an in-flight transaction drove it to
completion, as practical \mcas does~\citep{harris2002mcas, barnes1993method}.
Helping means several threads co-install one transaction's records, and
co-installation is exactly the condition of \cref{prop:abainstall}: a record can
be planted by a non-owner, so a stale late plant can resurrect a retired
descriptor. Closing that needed a per-record install latch---a
\code{FREE}/\code{BUSY}/\code{DONE} state machine on every record, arbitrating
which installer owned it.

The lesson is that this latch already made the helping engine bounded-blocking:
an installer that found a record \code{BUSY} waited on its holder. Helping never
bought lock-freedom; it bought a \emph{more complex} bounded-blocking engine.
And that complexity cost throughput in a way that can be counted before it is
measured. The latch is not merely state, it is \emph{atomics}, and---this is the
part worth being exact about---it does not replace the plant. A slot is claimed
by compare-and-swap under either design, because the latch guards \emph{that
record's} install and not the slot, which another transaction's record may also
target. Everything helping adds sits around an unchanged plant. Claiming the
record costs a compare-and-swap (\code{FREE}\,$\to$\,\code{BUSY}) and
releasing it costs a store. Settling it costs two more: the owner must first
claim the same word (\code{FREE}\,$\to$\,\code{DONE}), so that no late plant
can begin behind its back, and must then convert the slot by compare-and-swap
rather than by the release store a sole driver uses, because a helper may be
converting that slot at the same instant. Three extra compare-and-swaps per
record, then, against a sole driver's one to plant and one release store to
settle. The transaction's decision follows the same logic: the owner can no
longer publish its status with a release store, since it is not the only thread
that may write it, so the single most important store in the commit becomes a
compare-and-swap as well. Only then come the effects those atomics have---the
per-record state machine as a hot cache line contended by every co-installer, and
helper work that largely duplicated effort the owner would complete anyway. That diagnosis is not
only ours: recent work on \kcas reports the same failure, that under high
contention the helping which guarantees lock-freedom ``may cause excessive cache
invalidations and significant performance degradation''---and responds by
throttling the helpers rather than removing
them~\citep{unno2026helping}. \Cref{sec:varversion} takes up that branch.
Dropping the helpers
removes the co-installer, hence the hazard, hence the latch, hence the traffic,
leaving one driver per transaction and a single coarse settle-wait.

The removal is a measured win, and it is what retired the helping path. On a
three-structure skip-list move at 192 cores, the sole-driver engine sustains
$40.7$ against the helping engine's $33.4$ million key-moves per
second---$+22.0\%$---and the margin grows as the structures shrink and the
shared upper levels grow hotter, reaching roughly $+30\%$ on the smallest
structures measured; it also grows with the
writer count, from $+5\%$ at a single writer to $+22\%$ at 192, which is the
shape a contention argument should have.\footnote{\code{bench\_txn\_3skiplist},
width~3, jemalloc with a per-CPU arena, median of five 2\,s runs with the
address-space layout pinned, on a 2\,$\times$\,96-core EPYC 9654
(384 hardware threads). Run-to-run spread within a point is at most $1.1\%$, so
the separation is not noise; medians are reported rather than the best of a
batch. Both arms are built from engine commit \code{97443472} and differ in one
compile-time flag, the install discipline. The figure at 192 writers reproduced
across three independent capture sessions ($+22.1\%$, $+21.2\%$, $+22.0\%$);
the contended margin is quoted loosely because it did not, ranging over
$+27\%$ to $+30\%$ across those sessions while holding to $1.5\%$ within any
one of them. See \cref{sec:availability}.}

The comparison has to be made at that older commit, and it is worth saying why
rather than burying it. The helping install was \emph{deleted} from the engine
once it lost, so there is no build of the shipped tree that still contains it:
\code{97443472} is the last commit carrying both disciplines, which is what makes
a one-flag A/B possible at all. That leaves the obvious objection---that a
retired tree may not describe the shipped one---so we measured the sole-driver
arm on both. At the disclosed commit it is between $0.7\%$ and $2.4\%$
\emph{faster} than at \code{97443472}, at every writer count from 1 to 192. The
winner of the A/B is the engine that ships, slightly improved; nothing in the
intervening work moved it.

% ---------------------------------------------------------------------------
% NOTE (drafting; retained deliberately -- see README framing rule).
% RE-MEASURED 2026-07-28, then AGAIN with the address-space layout pinned.
% The layout matters: the sibling list benchmark is BIMODAL under ASLR (two
% clusters 19% apart), and while these skiplist panels are not visibly bimodal
% -- worst spread 1.1% -- the SERIALIZED lane arm is, which is why 64 writers is
% omitted from sec:escalation's table.  All four sweep scripts now pin it.
%
% WHAT REPRODUCES AND WHAT DOES NOT.  The +22% headline reproduced across three
% independent capture sessions (+22.1, +21.2, +22.0).  The CONTENDED margin did
% not: +27.0, +28.5, +29.9, and its peak moved from 1920 to 960 keys/skiplist,
% while within-run spread stayed ~1.5%.  So the text says "roughly +30%" rather
% than a two-decimal figure, deliberately -- do not tighten it back up without a
% fourth session that says the same thing twice.
%
% The stale 2026-07-10 numbers (39.9 vs 32.8, +21.7%,
% +25.8%) are gone; these are fresh.  Data + scripts live in
% efficios-trie-benchmark: scripts/p2_engine_remeasure.csv (240 runs, 0
% failures), scripts/build_p2_remeasure_arms.sh, run_p2_engine_remeasure.sh,
% analyze_p2_remeasure.py.  DUR=2000 RUNS=5, same 2x96c EPYC 9654, idle.
%
% WHY THE A/B IS AT 97443472 AND NOT AT THE PIN.  Helping was DELETED in
% a07b67ba, so it cannot be built at b3e23f9f at all.  97443472 is the last
% commit carrying both disciplines (-DURCU_MCAS_STOCK reverts to helping), so it
% is the only place a one-flag A/B exists.  The drift paragraph above is what
% answers the "but that tree is retired" objection: sole-driver measured at the
% pin too, and it is 0.5-3.1% FASTER there at every core count.  Do not delete
% that paragraph to save space -- without it the A/B describes a dead tree.
%
% THREE TRAPS FOUND WHILE PREPARING THE RE-RUN.  Each would have produced a
% plausible WRONG number rather than an error, and each is documented in
% run_p2_engine_remeasure.sh.  If these numbers are ever refreshed again, read
% that header first:
%   1. The committed engine A/B script still passes --ryw 1, a flag the bench
%      retired (1c1b124); the arg parser's `else usage()` exit(2)s, so re-running
%      it writes ZEROS into the CSV as if they were data.
%   2. At 97443472 URCU_TXN_RYW_DEFAULT is 0 and the bench's enable_ryw() call is
%      gone, so the historical arms need -DURCU_TXN_RYW_DEFAULT=1 or movesper 3
%      runs WITHOUT read-your-own-writes and corrupts instead of failing.
%   3. -DURCU_MCAS_STOCK flips FIVE things (install, AGE_ESCALATE, Bloom 16->1
%      words, Bloom k 3->1, AGE0_TRYLATCH).  The A/B must vary only the install,
%      so the helping arm restores the txn-layer defaults explicitly; otherwise
%      the "sole-driver win" silently absorbs the Bloom widening.  Arms are
%      verified distinct by counting uatomic_cmpxchg sites (7 vs 13).
%
% P1's lesson held: confirm the two arms do identical work apart from the
% mechanism under test.  Here that meant checking the 1-writer point, where
% optimistic and always-escalate come out identical (1.30032 vs 1.30030) -- an
% uncontended lane costs nothing, so the lane arm is measuring the funnel and
% not itself.
% ---------------------------------------------------------------------------
\subsection{Abort-others instead of waiting}
\label{sec:varabort}

A committer that meets a foreign transaction in its path could abort that
transaction rather than wait for it. This is obstruction-free contention
management, and it is well charted: DSTM's ``aggressive'' manager aborts any
conflicting transaction on sight~\citep{herlihy2003dstm}, and the priority and
timestamp policies that tame its livelock are Scherer and
Scott's~\citep{scherer2005contention}, whose root is the wound-wait rule of
classical database concurrency control~\citep{rosenkrantz1978woundwait}. This
engine had it, in the form of \emph{stealing}: a higher-ranked transaction evicted
a foreign one and took its slots. Stealing and helping were one subsystem---to
drive a foreign transaction is to help it, to evict one is to steal from it---so
they were removed together, and their combined cost is the measurement of
\cref{sec:varhelping}. The reason it loses is the same either way: aborting a
foreign transaction discards work that was often about to complete, and the retry
it forces costs more than the bounded wait that replaces it.

\subsection{Age-as-resolution}
\label{sec:varage}

A transaction could resolve a contended slot by \emph{version} rather than by
\emph{status}: stamp each transaction with an age, and let a reader or a
committer decide a slot by comparing ages instead of reading a status word. This
is the timestamp-ordering lineage---DSTM's timestamp
manager~\citep{herlihy2003dstm}, again over wound-wait---and it is a different
thing from the version-as-\emph{validity} of TL2 and
LSA~\citep{dice2006tl2, riegel2006lsa}, which we do not adopt and cite only to
distinguish. Age-as-resolution is rejected on a correctness trap, which is
decisive before throughput enters the argument: an age comparison cannot validate
a slot that a transaction reads but does not write, because there is no record to
carry the age, so the scheme cannot express the guarded read of
\cref{sec:structures} without reintroducing a status word anyway---at which point
it is the status mechanism with an extra field. We did not carry it to a separate
throughput measurement; the trap retired it first.

\subsection{DCAS or RDCSS for the install}
\label{sec:vardcas}

The classical fix for install-time A-B-A is a double-width primitive---DCAS, or
Harris, Fraser and Pratt's RDCSS~\citep{harris2002mcas}. \Cref{sec:norddcss}
showed sole-driver makes it unnecessary, but it was a live alternative before that
argument was found, and its costs are worth recording because they are the reason
it would have been a poor choice even where it is correct: DCAS is not portably
available and RDCSS adds a second conditioning word to every install, which is
hot-cache-line budget spent on a hazard the regime does not actually admit.

\subsection{Keeping the helpers, and paying in version bits}
\label{sec:varversion}

The alternative that most directly contradicts this paper's choice is to keep
helping and close its A-B-A hazard some other way. It is a live position, taken
independently and concurrently with this work: \citet{unno2026helping} retain
lock-free helping, regulate it with backoff and a helper count so that contention
does not collapse throughput, and suppress the install A-B-A by embedding a
\emph{version} in every target word, which lets a helper recognize an operation
it has already performed. Their diagnosis of helping under contention matches
ours (\cref{sec:varhelping}); the conclusion drawn is the opposite one.

We record it here rather than dispute it, because the two designs pay in
different currencies and the paper has not measured theirs. Version embedding
buys ABA-freedom under helping at the cost of bits taken from every transacted
word---a per-slot encoding budget this facility deliberately leaves to the
embedder (\cref{sec:constraints}), which spends it instead on a per-record tag it
chooses itself. And the trade only arises at all under helping: the hazard the
versions suppress is the one \cref{prop:abainstall} says a sole driver never
admits, so what is bought with encoding space there is had for nothing here. That
is the comparison worth making, and it is a design one; a throughput comparison
against an independent implementation belongs in the evaluation paper.

\subsection{Pure lock-freedom}
\label{sec:varlockfree}

Finally, one could insist on lock-freedom outright---no bounded wait, no
escalation lane. The measured history is what argues against it: lock-freedom was
unreachable the moment helping forced the per-record install latch
(\cref{sec:varhelping}), because the latch is itself a blocking point. Its one
genuine prize is immunity to preemption, and \cref{sec:preemption} argues that
prize is available more cheaply, for the common case, through preemption
\emph{avoidance} rather than a standing lock-free cost. We do not claim
lock-freedom is unachievable for this problem; we claim that the version of it we
could build lost, and that the reason it is attractive is answered by a cheaper
mechanism.

\begin{table}[t]
\centering
\footnotesize
\setlength{\tabcolsep}{5pt}
\begin{tabular}{@{}p{0.26\linewidth} p{0.40\linewidth} p{0.24\linewidth}@{}}
\toprule
\textbf{Alternative} & \textbf{Why rejected} & \textbf{Antecedent} \\
\midrule
Helping (co-install) &
  forced a per-record \code{FREE}/\code{BUSY}/\code{DONE} install latch $\Rightarrow$
  bounded-blocking anyway; hot-cache-line traffic; duplicated work. Removing help
  and steal: \textbf{$+22.0\%$} at 192 writers, to roughly $+30\%$ under contention &
  \citet{harris2002mcas}; \citet{barnes1993method} \\
Abort-others &
  the same subsystem as helping (steal $=$ evict a foreign txn); removed with it
  (\cref{sec:varhelping}); discards near-complete work &
  DSTM ``aggressive'' \citep{herlihy2003dstm}; \citet{scherer2005contention} $\leftarrow$ wound-wait \citep{rosenkrantz1978woundwait} \\
Age-as-resolution &
  cannot validate a read-not-written slot without a status word---a correctness
  trap, decisive before throughput &
  DSTM timestamp \citep{herlihy2003dstm}; \emph{not} TL2/LSA \citep{dice2006tl2, riegel2006lsa} \\
DCAS / RDCSS install &
  unnecessary under sole-driver (\cref{sec:norddcss}); unportable; hot-cache-line budget &
  RDCSS \citep{harris2002mcas} \\
Helping $+$ version-embedded words &
  not rejected on a number: an independent, concurrent design
  (\cref{sec:varversion}) that keeps helping and spends per-word encoding bits on
  the hazard a sole driver never admits &
  \citet{unno2026helping} \\
Pure lock-freedom &
  unreachable once co-install forced the latch; its one prize (preemption immunity) is cheaper via avoidance &
  --- \\
\bottomrule
\end{tabular}
\caption{The measured design space around the sole-driver facility. Every row but
the last was built or seriously considered here and then rejected; the last is
independent, concurrent work, recorded because it answers the same hazard the
opposite way and is not ours to reject. The decisive help-and-steal
removal ($+22.0\%$, \cref{sec:varhelping}) is the throughput evidence this paper
carries; the remaining rejected rows fall on portability or correctness grounds
that hold independently of a number, and the broader comparative study is
deferred. Disclosing each rejected approach is what makes the map defensive prior
art as well as an argument.}
\label{tab:variations}
\end{table}

% ===========================================================================
\section{Logical Deletion}
\label{sec:tombstone}

\Cref{sec:background} named a hazard the facility does not by itself close:
per-slot atomicity is not operation atomicity. A \kcas commits the slots it
names, all together; but a structural operation can be wrong because of a slot it
does \emph{not} name. The discipline that closes it is \emph{logical deletion}, a
known technique; this section states what it does and the design rationale that
fits it to this facility, and \cref{sec:structures} spends it on the worked
structures. It is presented here, at the facility level, because every mutating
structure on the facility needs it---not only the ones shown---and we claim no
novelty in the mechanism itself (\cref{sec:tombstoneclaim}).

\subsection{The hazard, concretely}
\label{sec:hazard}

Consider a doubly-linked list $A \leftrightarrow B \leftrightarrow C$, and two
operations racing: inserting a node after $A$, and deleting $B$. There are two
ways they can interact, and only one of them is closed by the facility's ordinary
slot contention.

The first is closed for free. Removing $B$ must rewrite $B$'s predecessor's
forward pointer---$A$'s \code{next}, the only forward pointer that names the path
through $B$---and an insert after $A$ rewrites $A$'s \code{next} too. So both
operations compare-and-swap the \emph{same} slot, \code{A->next}, with the same
expected old value, and the facility commits at most one of them. The loser fails
its old-value check, aborts, re-reads, and re-descends into a list that is still
coherent. Whichever wins, the outcome is a well-formed list. This race is closed
\emph{structurally}, by the fact that the two operations were contending for a
common slot after all---not by the tombstone.

The second is the real hazard, and it survives a perfect \kcas. An insert after
$A$ rewrites $B$'s \emph{back} pointer to name the new node, but it only
\emph{reads} $B$'s forward pointer---it has no reason to write it. A concurrent
deletion of $B$ writes $B$'s forward pointer and frees $B$. These are
\emph{different} words: the back-edge write and the forward-edge read do not
contend, so both the insert and the delete can commit, and the insert has just
written a back pointer into a node that is being freed. Nothing tore; the two
operations simply never met on a shared slot. Closing this needs two things: a
durable signal that $B$ is being deleted, set where a would-be attacher can check
it (\cref{sec:mark}), and a way to fold that check into the attacher's own commit
so its write is refused if $B$ died (\cref{sec:structures}).

\subsection{The tombstone is updater-only}
\label{sec:mark}

When a node is deleted, the operation that unlinks it also sets a
\emph{logical-deletion tombstone} on the node, as one word of the same multi-word
commit that unlinks it. Where the tombstone lives is a placement choice. It can
be packed into a reserved tag bit of the node's own forward pointer, distinct
from the facility's descriptor tag, which costs no extra word; or it can be a
separate status word, updated as an ordinary transaction record. Both are the
same logical object---a durable per-node ``this node is being deleted'' flag, set
atomically with the unlink and living for the node's grace-period lifetime---and
\cref{sec:tombstoneplacement} returns to why one might prefer either.

Its role is entirely among \emph{updaters}: once the tombstone is set, every
update transaction treats the node as \emph{frozen}, and the commit machinery
enforces the agreement---an update that would touch a frozen node fails to
commit. Two consequences follow, and they are the whole of what the tombstone is
for.

\emph{Deletion is single-updater.} A second delete of the node, or an insert or
delete anchored to it, finds the tombstone set and aborts, terminating with ``no
such entry'' rather than acting on a node that is already leaving. Exactly one
transaction ever logically deletes a node.

\emph{Nothing attaches to a departing node.} An operation that would link a fresh
element onto the node---writing one of the node's edges---must first agree the
node is still live, which it does by validating the tombstone into its own commit
(\cref{sec:guard}). If the tombstone is set, the attaching commit fails, so no
edge is ever installed into a node that is about to be reclaimed. This is the
disjoint-word hazard of \cref{sec:hazard} closed: the attaching write and the
tombstone are different words, but the validate makes the commit fail on the one
while it writes the other.

\subsection{Why readers never touch it, by design}
\label{sec:tombstoneplacement}

A reader consults the tombstone for nothing, and this is deliberate rather than
incidental. The freeze the tombstone announces is \emph{redundant} with the
transactions that enforce it: because every updater agrees not to touch a
tombstoned node, a reader that reached the node before its grace period sees an
object no one will mutate again---a frozen node---\emph{for free}, without loading
the tombstone at all. The updaters' agreement hands the reader exactly what a
per-read tombstone check could have told it, so the check would be redundant.

That the reader is off the tombstone is what makes its placement free, and it is
the series' read-side thesis carried one step past the companion paper. The
companion paper's proxy tag at least rode in the pointer a reader loads---free to
strip, but present in the reader's world. The tombstone need not be present at
all: packed into the forward pointer it is stripped for free on a word the reader
was loading anyway, and moved to a \emph{separate} word it leaves the reader's
forward pointer pristine and can sit on a cache-cold line the reader never
touches. Either way it adds nothing the reader must load, so in the steady state
the read side is byte-for-byte and cycle-for-cycle what it would be with no facility
present.

\subsection{Prior art}
\label{sec:tombstoneclaim}

% ---------------------------------------------------------------------------
% RESOLVED (Mathieu, 2026-07-18): no novelty is claimed for logical deletion.
% Disclose it as an applied known technique; cite Harris (packed marked pointer)
% / EFRvB / SCX finalize (separate-word flag). The one design property worth
% stating -- readers never consult the tombstone, so it can be cold-CL -- is
% DESIGN RATIONALE, not a claim; SCX is the closest prior art (its info field
% readers also never consult). Do NOT reintroduce a "composition" novelty claim,
% the "forwarding to old" framing (readers do not forward through the tombstone
% -- it is updater-only), or the GC-relocation (Brooks/ZGC/Shenandoah/Bronson)
% distinction apparatus, which existed only to defend the retracted forwarding
% claim. The SCOPE's DECISION #4 (claim the tombstone as a composition) is
% RETRACTED -- see the mark-not-proxy memory.
% ---------------------------------------------------------------------------

Logical deletion is old, in both of its placements, and we claim no novelty in it
as a mechanism---it enters this paper as one ingredient of the conjunction
(\cref{sec:relclaims}), not a result of its own. A marked pointer whose bit
logically deletes a node and fails a racing
writer's CAS is Harris~\citep{harris2001lists}; the equivalent as a separate flag
or record interpreted by updaters is the non-blocking-tree discipline of Ellen,
Fatourou, Ruppert and van Breugel~\citep{ellen2010nbbst} and Natarajan and
Mittal~\citep{natarajan2014bst}, and freezing a record set so later updates to it
fail is the \code{SCX} ``finalize'' of Brown, Ellen and
Ruppert~\citep{brown2013scx}. Setting the tombstone as one word of the same
multi-slot commit that unlinks the node---so that logical deletion and physical
unlinking share a single linearization point---is what a \kcas over a tree already
does with such a flag.

One property is worth calling out, as a design rationale rather than a claim: the
tombstone is consulted by updaters alone and never by readers. That is worth
separating from the technique it belongs to, because in the lineage the technique
comes from the mark is \emph{not} updater-only. A marked node in
Harris's list stays reachable until some thread physically unlinks it, so a
traversal that did not examine the mark would report a deleted node as present:
there the mark does double duty, coordinating updaters and telling readers the
node is gone. That duty is forced by the setting rather than chosen. Without a
grace period a node cannot be unlinked and freed while a reader may still hold
it, so logical deletion has to be a state the reader can see.

An \rcu base removes the second duty. Here the unlink \emph{is} the deletion as
far as a reader is concerned---the commit re-points the predecessor, the node
stops being reachable, and the grace period makes freeing it safe---so existence
is answered by reachability and never by a flag. What the tombstone adds is
addressed entirely to updaters (\cref{sec:mark}), and it reaches readers only
indirectly, through the updaters that decline to touch a frozen node. The nearest
prior art for a mark readers do not read is \code{SCX}, whose descriptor sits in
a field readers never consult; its reason is a linearization argument over plain
reads, where ours is the substrate. Keeping the tombstone off the reader's path
therefore costs nothing to arrange (\cref{sec:tombstoneplacement}) and leaves its
placement free---packed for compactness, or on a cold line for a pristine reader
pointer. We disclose that as the reason the technique fits this facility, and
claim nothing beyond the fit.

% ===========================================================================
\section{The Same Structures, Now Multi-Writer}
\label{sec:structures}

The companion paper showed a bidirectional list and a hash list on the
\fliplatch. The series keeps the same worked structures across papers, so that
what changes from one to the next is exactly the facility's capability, not the
example. Here they run under concurrent writers.

Every slot these two structures transact is a shared one: they are general-purpose
containers, so no lock is assumed over any of their links and every record they
record is of the shared kind. The exclusive kind of \cref{sec:kinds} is therefore
not exercised below---it is for the embedder who \emph{does} hold a lock over part
of what an edit touches, and who would otherwise have had to bridge two facilities
to move it atomically with the rest.

\subsection{Same-slot contention, handled by the facility}
\label{sec:contention}

The parts that are purely about two writers wanting the same edge need nothing
beyond \cref{sec:facility}. Adjacent inserts and deletes that re-point a common
\code{next} slot contend on it through descriptor-naming installs; the facility
picks a winner, the losers abort and re-descend, and the aging escalation of
\cref{sec:escalation} carries a persistently unlucky handle to a fair lane. A
composed operation that spans structures folds its edges into one commit through
the same read-your-own-writes discipline the companion paper introduced---which,
being shared by both facilities, is not this paper's to claim. There is no
list-specific machinery for contention; there is the facility.

\subsection{The one guarded read}
\label{sec:guard}

The disjoint-word hazard of \cref{sec:hazard} is what the structures must handle
themselves, and it needs exactly one new thing beyond the tombstone. An insert
after a node re-points that node's successor's \emph{back} pointer, but the slot
that serializes the insert against a deletion of the successor is the successor's
\emph{forward} pointer---a slot the insert reads but does not write. Because the
facility installs a transaction's records in slot-address order, it can reach the
back-pointer write before it reaches the shared forward slot, and so drive the
back-edge store against a successor a concurrent deletion is freeing.

The fix is to fold the read of that forward slot into the commit---to
\emph{validate the successor's tombstone}. The insert issues a
\emph{load-validate} on the successor's forward pointer, the slot where a deletion
of the successor sets its tombstone: it reads the value and records a load-only
guard---a record of the form \edit{v}{v}---so that the commit succeeds only if that
slot still resolves to $v$ at the install point. If a concurrent deletion has set
the tombstone in the meantime, the guard's old value no longer matches, the commit
fails, and the insert retries. The guarded read is the read-side companion to the
tombstone: the forward edge is made safe by a \emph{write}-validated slot, the back
edge by a \emph{read}-validated one.

A guard is of the shared kind necessarily, not by choice. What enforces it is the
old-value check of the install itself (\cref{sec:kinds}), so a guard is exactly a
shared record that happens to leave its slot unchanged; the exclusive install has
no check to fail and could not carry one. That is not a limitation: a slot the
embedder holds a lock over needs no guard, because the lock already excludes the
writer the guard would have caught.

% Inserting E between P and N under concurrent writers, and the same insert
% aborting.  The companion paper's fig:listinsert adapted to the multi-writer
% case, which changes what the figure is FOR.
%
% P1's version is a four-phase lifecycle (build | install | commit | settle),
% because under exclusion the interesting thing is that a commit cannot fail and
% the reader sees one coherent list either side of it.  Here the commit CAN fail,
% so the interesting thing is the fork: one write set, one store, two outcomes.
% Drawing the two outcomes as separate figures would duplicate the shared write
% set and split the comparison across a page break; forking shows that the
% install is common and only the terminal value differs, which is the point.
%
% Two things this must make obvious that no other figure in the paper does:
%   (a) the GUARD -- a record whose old and new are the SAME value, read but not
%       written (sec:guard).  A table row with v in both columns explains
%       load-validate faster than the paragraph does.
%   (b) an abort is not a partial edit.  The FAILED branch leaves the list
%       byte-for-byte as it was found, which is what makes a losing writer's
%       retry safe.
%
% Deliberately NOT redrawn here: the slot -> record -> status anatomy, which is
% fig:tristate's job.  This figure starts from the write set and assumes it.
%
% Layout note: render the page and LOOK at it.  Every fault in this paper's other
% two figures passed a clean build.
\begin{figure}[t]
\centering
\begin{tikzpicture}[
  font=\small,
  lnode/.style   = {draw, rounded corners=2pt, minimum width=8mm, minimum height=5.5mm, fill=gray!8},
  nnode/.style   = {draw, rounded corners=2pt, minimum width=8mm, minimum height=5.5mm, fill=green!14},
  sel/.style     = {draw, very thick, minimum width=15mm, minimum height=8mm, fill=orange!18},
  trow/.style    = {anchor=west, inner sep=0pt, minimum height=6mm, font=\footnotesize},
  biarr/.style   = {latex-latex, semithick},
  arr/.style     = {-latex, thin},
  phase/.style   = {font=\footnotesize\itshape, text=black!70},
]

% ===== the write set: two writes and one guard ==============================
\node[trow, fill=blue!14] (hdr)  at (0, 3.30)
  {\makebox[15mm][c]{\textit{slot}}\makebox[12mm][c]{\textit{old}}\makebox[12mm][c]{\textit{new}}};
\node[trow, fill=blue!6]  (rec1) at (0, 2.70)
  {\makebox[15mm][c]{\code{P.next}}\makebox[12mm][c]{\code{N}}\makebox[12mm][c]{\code{E}}};
\node[trow, fill=blue!6]  (rec2) at (0, 2.10)
  {\makebox[15mm][c]{\code{N.prev}}\makebox[12mm][c]{\code{P}}\makebox[12mm][c]{\code{E}}};
\node[trow, fill=orange!10] (rec3) at (0, 1.50)
  {\makebox[15mm][c]{\code{N.next}}\makebox[12mm][c]{$v$}\makebox[12mm][c]{$v$}};
\draw[thin, blue!45] (0, 1.2) rectangle (3.9, 3.6);
\draw[thin, blue!45] (1.5, 1.2) -- (1.5, 3.6);
\draw[thin, blue!45] (2.7, 1.2) -- (2.7, 3.6);
\draw[thin, blue!45] (0, 3.0) -- (3.9, 3.0);
\draw[thin, blue!45] (0, 1.8) -- (3.9, 1.8);

% Under the table, not beside it: to the right of rec3 is where the FAILED
% branch runs, and the note printed straight over it.
\node[font=\scriptsize, text=orange!45!black] at (1.95, 0.8)
  {\textbf{guard}: read, never written};
\node[above=0.5mm of hdr, font=\footnotesize\itshape] {the write set};

% ===== the status word, forking =============================================
% Elbow arrows, not diagonals: the two branches leave the status word vertically
% and run out horizontally, which keeps their labels on a clear vertical segment
% instead of riding a slope through the outcome text.
\node[sel] (st) at (7.0, 2.4) {\code{status}};
\draw[arr] (3.9, 2.4) -- node[above, font=\scriptsize] {install} (st.west);

\draw[-latex, semithick] (7.0, 2.8) -- (7.0, 4.0) -- (9.15, 4.0);
\node[right, font=\scriptsize] at (7.05, 3.45) {\SUCCEEDED};

\draw[-latex, semithick, gray!60!black] (7.0, 2.0) -- (7.0, 0.6) -- (10.0, 0.6);
\node[right, font=\scriptsize, text=gray!45!black] at (7.05, 1.35) {\FAILED};

% ===== outcome 1: committed =================================================
% Node centres 1.7 apart against an 8mm box, so the link arrow gets a 9mm span.
% At the 1.1 spacing this started from, the gap was 3mm and two arrowheads ate
% all of it -- the boxes read as touching. Both outcome rows and all three notes
% below them share the centre 11.4, so the rows stay stacked.
\node[lnode] (cP) at ( 9.7, 4.0) {\code{P}};
\node[nnode] (cE) at (11.4, 4.0) {\code{E}};
\node[lnode] (cN) at (13.1, 4.0) {\code{N}};
\draw[biarr] (cP) -- (cE);
\draw[biarr] (cE) -- (cN);
\node[font=\scriptsize, text=green!35!black] at (11.4, 3.35)
  {\code{E} linked, both directions};

% ===== outcome 2: aborted ===================================================
\node[lnode] (aP) at (10.55, 0.6) {\code{P}};
\node[lnode] (aN) at (12.25, 0.6) {\code{N}};
\draw[biarr] (aP) -- (aN);
\node[font=\scriptsize, text=gray!45!black] at (11.4, -0.05)
  {unchanged: every record resolves to \textit{old}};
\node[font=\scriptsize, text=gray!55!black, align=center] at (11.4, -0.8)
  {a concurrent delete of \code{N} marked it,\\so the guard's \textit{old} no longer holds};

\end{tikzpicture}
\caption{One write set, one store, two outcomes---the companion paper's insert
under concurrent writers. Inserting \code{E} between \code{P} and \code{N} moves
two edges, and needs one thing more that the single-writer case did not: a
\emph{guard} on \code{N}'s forward pointer, a record whose old and new are the
same value $v$. No other record in the transaction names that slot---the insert
reads it but never writes it---so without the guard the back-edge record could
install into \code{N} before the transaction discovered, at \code{P}'s forward
pointer, that it had to abort (\cref{sec:guard}). On \SUCCEEDED\ every record resolves to its new value and
\code{E} is linked in both directions at once. On \FAILED---here because the
guard's expected value no longer holds---every record resolves to its old value
instead, so the list is left exactly as it was found and the writer may simply
re-descend and retry. An abort is never a partial edit.}
\label{fig:listinsert}
\end{figure}


This is deliberately the minimal, functional form of a general facility. The
guard exists here for one reason---to make this paper's own worked example
correct under concurrent writers---and it is presented as exactly that: a single
read folded into a single commit's conflict set. The general read-side model it
is an instance of---a taxonomy of read policies, à-la-carte guards on arbitrary
read-not-written slots, and the declarations that tune them---is the subject of
the next paper, and this paper claims none of that generality. It claims only
that its own list is usable, and shows the one guard that makes it so.

% -----------------------------------------------------------------------
% NOTE (drafting; retained deliberately -- see SCOPE DECISION #5).
% This is the P2/P3 boundary. P2 introduces ONE guard, functionally, as the
% back-edge dual of the mark, to make its OWN flagship example correct under
% MW. It must NOT claim the a-la-carte read-set-validation novelty (that is
% P3-1, whose novelty is the barrier-free-RCU substrate + opt-in-per-slot
% direction). Keep this subsection about "our list needs it", never about
% "here is a general facility". The exact source phrasing is that the
% load-validate of succ->next "exists ONLY to guard the pprev write".
% -----------------------------------------------------------------------

\subsection{Why the hash list transacts its back link, when the single writer
did not}
\label{sec:hlistpprev}

The companion paper's hash list wrote its backward link---each node's
\code{pprev}, a pointer to the slot that names it---with a plain store, on the
grounds that no reader ever loads it and a single updater has no peer to race.
Under concurrent writers that reasoning does not transfer, and the reason is
instructive about the whole difference between the two disciplines.

For the single writer, ``writer-only'' meant ``private'': a plain store to
\code{pprev} was instantly visible to the sole updater and to no one else, so
there was no window in which it could be seen wrong. Under concurrency
``writer-only'' still means shared: peer mutators read \code{pprev} while
resolving their own edits, and it is lifetime-critical. A plain store cannot be
made safe either eagerly or lazily. Store it before the commit and a lost commit
leaves the store behind, so a live node names a slot inside a node that never
linked. Store it after the commit and there is a window in which \code{pprev} is
not a stale hint but a dangling pointer into another node---and such hints
accumulate on a single slot, unbounded, uncovered by any read-side critical
section because they sit in live nodes indefinitely. So under concurrency
\code{pprev} is transacted: the multi-writer equivalent of the single writer's
``eager'' store is ``atomic with the commit,'' which is the transaction. This is
not a regression; it is the exclusion crutch coming due, in one concrete slot.

It is worth asking whether \cref{sec:kinds} rescues the plain store after all, since
an exclusive record \emph{is} a plain store inside the commit. It does not, and
the reason shows the contract has teeth. Exclusive is a promise that the embedder
holds a lock excluding \emph{every} writer of that slot, and a general-purpose
hash list makes no such promise: \code{pprev} is written by whichever peer mutator
is re-linking around this node, under no lock at all. Recording it as exclusive
would be exactly the globally-inconsistent kind of \cref{sec:constraints}, and the
plain park would race the very peers the transaction exists to serialize against.
The kinds let an embedder \emph{spend} exclusion it already has; they do not
manufacture it.

\subsection{Reclamation: no re-link before a grace period}
\label{sec:reclaim}

The build--publish--reclaim discipline carries over, with one multi-writer
addition that the single-writer paper explicitly did not have. A deleted node
must not be re-linked---re-inserted, or recycled into a fresh node---before a
grace period, and not only for the usual reason that a reader may still hold it.
The deleting transaction may still hold a parked descriptor in that node's
forward slot, awaiting its owner's settle. An insert builds its new node's links
with plain stores while the node is private; a plain store into a slot that still
holds a parked descriptor is overwritten when that descriptor settles, and the
write is silently lost, leaving a permanently incoherent edge. So a node's
storage must wait out a grace period before reuse exactly as its reclamation
would---the same wait, extended from ``free'' to ``reuse.''

% ===========================================================================
\section{Limitations}
\label{sec:limitations}

The facility's limits divide the same way the companion paper's did: some are
properties of the model, which can be disliked but not violated, and the rest are
contracts on the embedder, which can be violated at a cost.

\paragraph{Per-slot atomicity is not operation atomicity.} A \kcas makes its
named slots change together; it does not make an operation correct with respect
to a slot it does not name. The disjoint-word structural races this admits are
real, and \cref{sec:tombstone,sec:structures} are the discipline that closes them
for the shown structures---the logical-deletion tombstone and the one guarded
read. A different structure with a different coupling owes itself the same kind of
analysis; logical deletion is a general technique, but recognizing which reads a
given operation must guard is the embedder's obligation.

\paragraph{Bounded-blocking, not lock-free.} The facility waits on a foreign owner,
so a preempted owner blocks its waiters until the escalation resolves them
(\cref{sec:boundedblocking}). It is starvation-free under a fair scheduler, not
lock-free, and correctness is independent of the preemption-avoidance lever of
\cref{sec:preemption}, which is not present in the measured engine and rescues
the tail rather than the correctness.

\paragraph{The slot's kind is unchecked, and it is global.} Of the embedder's
contracts this is the one with the worst failure mode, because it is the only one
the facility cannot see. Naming a slot exclusive asserts something about
\emph{every other thread}---that the embedder's lock excludes all of them from it
---and a transaction sees only its own records, so a disagreement between two
threads is invisible until it corrupts (\cref{sec:constraints}). Nothing here
detects it, and no amount of testing establishes it: it is discharged by the
embedder's locking design or not at all. The safe direction is free, which is the
mitigation on offer---naming a slot shared is correct unconditionally, so the
contract is only ever at risk where an embedder reaches for the optimization.

\paragraph{Value-CAS atomicity.} A record is validated to hold its old value at
the linearization point, not to have been stable throughout the transaction. An
operation whose correctness requires detecting an away-and-back change must carry
its own version in the word; snapshot and version semantics are layered on top,
as with any value-based \mcas. Memory safety holds unconditionally; atomicity is
conditional on the embedder's contracts---the per-record tag kept free,
pairwise-distinct slots, deferred descriptor reclaim, and the logical-deletion
discipline where
operations couple through unwritten slots.

% ===========================================================================
\section{Related Work}
\label{sec:related}

\subsection{Multi-word compare-and-swap}

The facility's descriptor-and-status shape
is that of practical \mcas~\citep{harris2002mcas, fraser2004lockfreedom,
fraser2007locks}, and its no-RDCSS result is stated against that line: RDCSS is
the machinery~\citep{harris2002mcas} that a recent \kcas
analysis~\citep{guerraoui2020mcas} ties to helping, which is what
\cref{sec:norddcss} removes. Descriptor reuse and lifetime are Arbel-Raviv and
Brown's~\citep{arbelraviv2017descriptors}, whose \rcu reclamation baseline over
this library is the prior art \cref{sec:settle} is careful not to reclaim as its
own.

\subsection{Blocking that beats lock-free}

That a blocking design can outperform
a lock-free one at scale is not a new observation---it is the lesson of broad
synchronization studies~\citep{david2013everything} and of flat
combining~\citep{hendler2010flatcombining}, which serializes deliberately to win
throughput. We claim the specific engine and its measured falsification trail, not
the general insight; the progress taxonomy we place it in is
standard~\citep{herlihy2003obstruction, herlihy2011progress}, and the fair-lane
fallback is an MCS lock~\citep{mellorcrummey1991mcs}.

\subsection{Mixing synchronization disciplines within one transaction}
\label{sec:relmixed}

The design point of \cref{sec:kinds}---each datum governed by one discipline,
one transaction spanning several---is not new, and it is charted most clearly in
the database literature. Tang and Elmore enumerate the four ways concurrency
control can be mixed and name this one exactly: \emph{one protocol per record,
multiple protocols per transaction}, with the protocol a global property of the
record and a transaction free to touch records under
several~\citep{tang2018cormcc}; \citet{wang2016mocc} is a prior instantiation,
mixing optimistic control with two-phase locking. We claim nothing about the
design point. Our contract that a slot's kind is global across all writers
(\cref{sec:constraints}) is the same discipline their framework imposes when it
requires every candidate protocol to be strict, and for the same reason.

What differs is what the mixing costs, and that is a property of the setting
rather than of the idea. \Citet{tang2018cormcc} define the overhead of mixing as
executing more than one protocol per transaction \emph{plus} synchronizing
concurrent accesses across protocols on each record, and their framework earns
the second down to nothing by constraining which protocols may be combined; its
commit executes each protocol's commit phase in turn, so atomicity across
protocols is established there by a serializability argument over separate
commits rather than by a single instant.
% check-terms will flag "serializability" on the line above. It is DELIBERATE and
% it describes CormCC, not this facility -- the subject is "its commit", theirs.
% The contrast being drawn is precisely that we do NOT need such an argument
% (one status store is the linearization point). Do not let an editing pass
% attach the word to our own mechanism; see the README framing rule.
Here both
terms vanish for a narrower reason. The two disciplines are not two protocols
that must be reconciled but one resolve path with two ways of writing into it, so
there is nothing to synchronize across them and no argument to make: the one
status store is the linearization point for both kinds at once
(\cref{prop:kindinvisible}). The first term goes the way of any special case,
by detecting that a write set does not actually mix. The narrowness is the point:
that reduction is available because both disciplines were built to share a
descriptor, and it says nothing about mixing protocols that were not.

\subsection{Logical deletion}

The tombstone of \cref{sec:tombstone} is applied
prior art, claimed as none of ours. Marked-pointer deletion is
Harris~\citep{harris2001lists}; the flag/record discipline in non-blocking trees
is Ellen et al.~\citep{ellen2010nbbst} and Natarajan and
Mittal~\citep{natarajan2014bst}; freezing a record set against later updates is
\code{SCX}~\citep{brown2013scx}, which is also the nearest prior art for the one
design property we call out---that readers never consult the tombstone
(\cref{sec:tombstoneclaim}).

\subsection{RCU trees, and reader-visible old versions}

Concurrent updates to an
\rcu search tree with per-node-locked writers are Arbel and Attiya's
Citrus~\citep{arbel2014citrus}; the sole-driver facility replaces those locks with
lock-free \mcas writers over the same \rcu reader discipline. Reader-visible old
versions under concurrent writers are RLU~\citep{matveev2015rlu} and
MV-RLU~\citep{kim2019mvrlu}, which give a coherent snapshot the \pseudotxn model
declines; a throughput comparison against them crosses guarantee classes and is
the evaluation paper's to make carefully. Transactional
boosting~\citep{herlihy2008boosting} is the composition pattern that pairing a
lock with this facility converges on.

\subsection{What this paper claims}
\label{sec:relclaims}

The paper makes two claims, at two scopes.

\textbf{The first is independent and specific: a sole-driver \mcas needs no
RDCSS.} The install-time A-B-A that RDCSS exists to prevent is helper-induced
(\cref{prop:abainstall})---it needs a second, foreign driver of one's own
records---so with no helper it cannot arise, and a plain value-CAS install is
A-B-A-safe by construction rather than by a double-word primitive. Value-CAS
install, single-driver ownership~\citep{israeli1994disjoint}, and the tie between
RDCSS and helping~\citep{guerraoui2020mcas} are each old; what we are not aware of
drawn this way is the conclusion---that dropping helping makes RDCSS
\emph{unnecessary}, not merely cheaper. It is a negative result, and it stands on
its own.

\textbf{The second is the conjunction, in the manner of the companion paper.} Each
remaining ingredient of the facility is separately old: the descriptor-and-status
shape is Harris's~\citep{harris2002mcas, fraser2004lockfreedom}; that a
\emph{blocking} design can beat a lock-free one is not our
insight~\citep{david2013everything, hendler2010flatcombining}; the progress
classes we place it in are standard~\citep{herlihy2003obstruction,
herlihy2011progress} and its fair serial lane is an MCS
lock~\citep{mellorcrummey1991mcs}; reclaiming descriptors after a grace period is
\rcu's ordinary treatment~\citep{arbelraviv2017descriptors}; logical deletion
is Harris's marked pointer and its descendants~\citep{harris2001lists,
ellen2010nbbst, brown2013scx}; and governing each datum by one discipline while
letting a transaction span several is charted in the mixed concurrency control of
\citet{tang2018cormcc} and instantiated before us by \citet{wang2016mocc}. What
we have not found is their assembly into one
facility: RDCSS-free and sole-driven, \textbf{bounded-blocking with a cost-scaled
adaptive coarseness}, \textbf{refcount-free to reclaim} at the instant owner-only
settle makes a transaction unreachable, \textbf{carrying both exclusion
disciplines at once} so that a commit spans locked and unlocked slots against one
linearization point, and \textbf{usable in place} through a
logical-deletion discipline that charges the reader nothing. The first claim is
the linchpin of this second one.

The fourth of those deserves its boundary drawn, since \cref{sec:relmixed} grants
the design point to others. What is ours is not that disciplines can be mixed but
that mixing these two costs nothing on either side: the kind never reaches a
reader (\cref{prop:kindinvisible}) and a write set that does not mix does not pay
(\cref{sec:kinds}). That falls out of a representational choice---one resolve
header shared by both kinds---which dissolves what had looked like an
incompatibility between the two facilities rather than defending against it. It is
the same shape of result as the first claim, a mechanism turning out to be
unnecessary; it is weaker because no one else had stated the obstacle, and we
claim it only as an element of this conjunction.

Everything else is disclosed, not claimed---and for a defensive publication that
is the point, not a shortfall: we describe the facility, and the alternatives we
rejected (\cref{sec:variations}), in enough detail to be rebuilt from these pages
alone.

% ===========================================================================
\section{Conclusion}
\label{sec:conclusion}

The companion paper gave the \rcu writer a second atomic operation under exclusion. This
paper removes the exclusion and pays for it in machinery the first had no need
of---a descriptor-naming install, a status word, a bounded wait, an escalation
lane---while refusing to pay for one thing that machinery usually drags in.
Because a transaction is driven only by its owner, the install-time A-B-A that
RDCSS exists to prevent cannot arise, so the facility installs with a bare
value-CAS and needs no RDCSS. That is a negative result---the first of the paper's
two claims, and the one it turns on.

The exclusion turns out to be droppable one slot at a time rather than all at
once. A record carries a kind, so a slot with concurrent writers is planted with
a CAS and a slot the embedder already holds a lock over is parked with a plain
store, both in the same commit and both decided by the same status store. The two
were separate facilities only because their proxies were separate layouts; giving
them one resolve header makes the difference writer-side and invisible to a
reader, which is what lets an embedder pay the blocking-progress price on exactly
the slots that are shared and on no others.

The facility is bounded-blocking, not lock-free, and we have argued that this is
the right target rather than a shortfall. The argument is measured: the
lock-free-adjacent alternative, helping, was built, and it forced a per-record
install latch that made it bounded-blocking anyway and cost it the throughput it
was supposed to buy. Removing the helpers removed the hazard, the latch, and the
traffic at once. What waits now is a single coarse settle-wait, escalating under
pressure into a FIFO-fair lane whose budget scales with an operation's cost, so
that cheap operations serialize early and expensive traversals stay parallel.

Making a concurrent structure usable took one known technique beyond a wide
commit, because per-slot atomicity is not operation atomicity. A deleted node
carries a logical-deletion tombstone, set atomically with its unlink, that update
transactions treat as freezing the node: deletion stays single-updater, and
nothing attaches to a departing node. Because that freeze is updater-only, a
reader gets a frozen node for free without ever consulting the tombstone, so it
stays off the reader's path and can sit on a cold cache line---the read side
unchanged, as the companion paper promised. Its read-side companion is a single
guarded read that folds one liveness precondition into the commit. The technique
is not ours; folded into the facility it is one part of the paper's second claim.
With it, the same bidirectional list and hash list the companion paper used are
usable under concurrent writers---the forward edge made coherent by a shared-slot
write, the back edge by a read-validated one, both coherent at every step.

That second claim is broader, and weaker in kind: each part---the
descriptor-and-status mechanism, the bounded-blocking progress, the refcount-free
reclamation, the logical deletion---is borrowed, but their assembly into one
facility, with the no-RDCSS result as its linchpin, is a combination we have not
found before.

\paragraph{The rest of the series.} The read-side programming model this paper
forward-references---the read-policy taxonomy, the general guard facility, and the
conflict and aging declarations---is the next paper's, demonstrated on the same
structures. The comprehensive comparative evaluation, against existence, RLU, and
MV-RLU, and across scaling and tail latency, is later still: it does little
defensively, which is why it is last, and it is the strongest academic artifact,
which is why it is worth doing carefully rather than early.

% ===========================================================================
\section*{Availability}
\label{sec:availability}

% ---------------------------------------------------------------------------
% Pin the COMMIT, never the branch (see README). b3e23f9f is the urcu-txn-dev tip
% carrying BOTH things this paper's mechanism sections depend on: the read-policy
% "re-verb from helping to waiting" (7be75ff6, load-bearing for the sole-driver
% read story, sec:soledriver) and the mixed sw/mw front-end (2a599d77..610d4792,
% load-bearing for sec:kinds). P1 pins the same commit. Verified pushed to
% github-dev/urcu-txn-dev.
% NOTE: the old comment here still named c7f428a8, a pin two revisions stale --
% the hash in the PROSE below was updated without it. Keep them in step.
% ---------------------------------------------------------------------------
The engine described here is free software, distributed under
\code{LGPL-2.1-or-later} as part of Userspace \rcu. Every mechanism this paper
describes is present in commit
\code{b3e23f9f558ce895aebe300ddb33aee9281137dc} (2026-07-25), on branch
\code{urcu-txn-dev} at \url{https://github.com/compudj/userspace-rcu-dev}. The
commit hash, not the branch, identifies the disclosed version: the branch is a
development tip and will move.

Within that tree, the material of this paper is the following headers, all under
\code{include/urcu/}:

\begin{description}\setlength{\itemsep}{2pt}
\item[\code{rcu-txn-mcas.h}] The sole-driver \mcas engine: records, the tri-state
      status word, descriptor-naming installs, the owner-only drive and settle,
      and the no-RDCSS argument in its own header block
      (\cref{sec:facility,sec:norddcss,sec:settle}). It carries both install
      disciplines of \cref{sec:kinds} behind one descriptor and one status word:
      the record's \code{kind} selects the CAS plant described here or the
      exclusive park of the single-writer paper, the two commit paths are
      \code{urcu\_txn\_desc\_commit()} and its \code{\codeunderscore sw} variant,
      and the resolve path is common to both.
\item[\code{rcu-txn.h}] The transaction front-end, and at the pinned commit the
      canonical one: it is the mixed front-end, so \code{urcu\_txn\_store\_mw()}
      and \code{urcu\_txn\_store\_sw()} record the two kinds into a single
      transaction. It also carries the read policy, read-your-own-writes,
      the disjoint and expect-conflict declarations, and the aging-to-fair-mutex
      escalation whose budget scales with operation cost
      (\cref{sec:escalation}). At this commit there is no separate
      multi-writer-only front-end: a write set of a single kind is a special case
      of this one, not a different facility.
\item[\code{rcu-txn-list.h}, \code{rcu-txn-hlist.h}] The multi-writer
      bidirectional list and hash list: the logical-deletion tombstone, the one
      load-validate guard, and the transacted back link
      (\cref{sec:tombstone,sec:structures}).
\item[\code{rcu-txn-status.h}] The commit-status enum both front-ends return, so
      an embedder tests a commit the same way whichever facility it drives.
\end{description}

\noindent The fair mutex of \cref{sec:escalation} is \code{fair-mutex.h}, and the
read-your-own-writes lookup filter both facilities share is \code{rcu-txn-bloom.h};
the single-writer \fliplatch of the companion paper is \code{rcu-txn-sw.h} and its
structures. The \code{rseq} time-slice-extension lever of \cref{sec:preemption} is
\emph{not} part of this engine and is named there as a design direction, not a
component of the disclosed tree.

This paper is licensed CC BY 4.0. It is published to place the design, and the
alternatives of \cref{sec:variations}, in the public record as prior art, so that
they remain free for anyone to implement and use.

\bibliographystyle{plainnat}
\bibliography{../common/urcu-txn}

\end{document}
